\documentclass[10pt]{article}
\usepackage{amsmath, amssymb}
\usepackage[round,authoryear]{natbib}
\usepackage{hyperref}
\usepackage{xcolor}
\usepackage{appendix}
\usepackage{booktabs}
\usepackage{graphicx}
\usepackage{caption}
\usepackage{float}
\usepackage{dsfont}
\usepackage{algorithm}
\usepackage{algpseudocode}
\usepackage{mathptmx}
\usepackage[T1]{fontenc}
\usepackage[margin=1.25in]{geometry}
\usepackage{setspace}
\setstretch{1.25}
\usepackage{threeparttable}
\captionsetup{labelfont=bf, textfont=it}
\usepackage{array}

% Configure hyperref for colored citations
\hypersetup{
    colorlinks=true,    % Color links instead of boxes
    linkcolor=blue,     % Color of internal links
    citecolor=blue,     % Color of citation links
    urlcolor=blue       % Color of URL links
}

\renewcommand{\appendixname}{Appendix}

\author{Bruno Cittolin Smaniotto\thanks{Replication material available at \url{https://github.com/brunosmaniotto/An-Adaptive-Hybrid-New-Keynesian-Model}.}}
\date{}

\begin{document}

\title{An Adaptive-Hybrid New-Keynesian Model}

\maketitle

\begin{abstract}
We develop an adaptive model of inflation expectations in which agents dynamically choose between anchoring to the central bank's target and extrapolating from recent inflation based on forecasting performance. When inflation is stable, anchored forecasts outperform and credibility remains high; when inflation deviates persistently, agents de-anchor and inflation becomes self-reinforcing. The framework provides a dynamic theory of central bank credibility, measured by the fraction of agents who trust the target, in which credibility and inflation outcomes evolve together through mutual feedback. This perspective sheds light on the 'persistence puzzle' and the apparent instability of Phillips curve estimates: both reflect time-varying credibility rather than structural breaks. Simulations show the framework is consistent with key features of the Volcker disinflation, the missing disinflation of 2008--09, and the post-pandemic inflation surge. Extensions incorporating trend-following expectations address hyperinflation dynamics and Brazil's Real Plan, while experience-weighted memory speaks to Japan's deflation trap and the persistence of inflation scars across generations. The framework implies state-dependent optimal policy: high-credibility central banks can `look through' temporary shocks, while low-credibility banks must respond aggressively.
\end{abstract}

\newpage

\section{Introduction}

Expectations play a central role in modern inflation theory. They determine how quickly inflation responds to shocks, how long policy takes to work, and how much disinflation costs. Despite their importance, the empirical record presents a number of puzzles that existing models struggle to reconcile. Why does the Phillips Curve appear steep during the 1970s but flat during the Great Moderation? Why did the Great Recession produce so little deflation despite massive slack? And how did inflation fall alongside unemployment in 2023-24? 

At the heart of these puzzles is the question of how inflation expectations are formed, an issue on which mainstream models have swung between extremes. Before the Lucas critique, models assumed backward-looking expectations, largely setting aside agents' ability to anticipate policy. The subsequent methodological shift was to assume Full-Information Rational Expectations (FIRE), which raises the opposite concern: agents become implausibly sophisticated and well-informed. A growing body of evidence from behavioral economics and survey data suggests the truth lies somewhere in between. Agents are neither naive extrapolators nor fully rational forecasters, but adaptive learners who adjust their forecasting strategies based on experience and the economic environment.

In this paper, we propose a model in this adaptive tradition, where expectation formation is shaped by both cognitive constraints and the economic environment. More specifically, agents dynamically choose between anchored and backward-looking forecasts based on their recent predictive accuracy, while retaining a natural cognitive preference for the simpler, static target when both rules perform similarly, following a learning rule inspired by reinforcement learning algorithms, with central bank credibility emerging endogenously from the resulting dynamics. The learning process requires minimal cognitive demands: agents perform no reasoning about the economy's structure or other agents' beliefs. They simply compare forecast errors using publicly available information, the announced target and realized inflation, and gravitate toward whichever rule has worked better. This framework can explain two seemingly contradictory facts. Inflation responds only modestly to shocks when it is low and stable \citep{coibion2015, ball2011}, yet becomes difficult to tame once it rises persistently \citep{erceg2003imperfect, goodfriend2005}.

Beyond this core mechanism, the framework sheds light on several phenomena that standard models struggle to explain, including why inflation persistence varies substantially across monetary regimes, why the Phillips curve appears steep in some eras and flat in others, and why comparable shocks, such as the oil episodes of the 1970s and 2008, can produce very different outcomes depending on prevailing credibility. The model provides a structural interpretation for the findings of \citet{hazell2022slope} and \citet{coibion2015}, who document that anchored expectations are central to understanding inflation dynamics. In our framework, anchoring is not assumed but earned, and its variation across regimes helps explain the apparent instability of reduced-form relationships.

We provide evidence that U.S. inflation persistence varies substantially across monetary regimes, consistent with the model's central prediction. Simulations show the framework captures key features of the Volcker disinflation, the missing disinflation of 2008–09, and the post-pandemic inflation surge. Extensions incorporating trend-following expectations address hyperinflation dynamics and stabilization episodes like Brazil's Real Plan, while experience-weighted memory speaks to Japan's deflation trap. The underlying logic is a ``gear-shift'' in which agents adopt progressively more complex heuristics as simpler ones fail, and these cognitive shifts feed back into inflation dynamics. The framework also implies state-dependent optimal policy: high-credibility central banks can "look through" temporary shocks, while low-credibility banks must respond aggressively.

\paragraph{Related Literature} The puzzles motivating this paper have been highlighted by recent empirical work. \citet{hazell2022slope} document that the Phillips curve is remarkably flat and attribute the apparent variation in slope across eras to expectation anchoring, but do not model how anchoring evolves. Their recent work on monetary policy \citep{nakamura2025taylor} further emphasizes that standard models struggle to match the observed transmission of policy to inflation. \citet{coibion2015} highlight the central role of expectations in resolving the ``missing disinflation'' of the Great Recession, while \citet{coibion2025inflation} argue that the post-COVID inflation surge reflects de-anchored expectations combined with supply shocks. \citet{weber2025} provide experimental evidence that attention to inflation is state-dependent: agents are inattentive when inflation is low but become highly attentive precisely when policy appears to be failing. These papers establish that expectation anchoring matters and that it varies with the environment; our contribution is to help explain how and why.

Our framework draws on two literatures. First, we build on behavioral macroeconomics, particularly work showing that bounded rationality has first-order macroeconomic consequences. \citet{gabaix2020} shows how cognitive limitations reshape the New Keynesian model, \citet{garciaschmidt2019} analyze how finite reflective depth affects expectation formation, and \citet{angeletos2023confidence} demonstrate that confidence dynamics alter the propagation of demand shocks. Second, we connect to work applying machine learning to model human cognition. \citet{binz2023} demonstrate that reinforcement learning algorithms accurately describe human decision-making in cognitive tasks, and \citet{charpentier2021} survey applications of these methods in economics. The specific reinforcement learning framework we adopt, the Multi-Armed Bandit, is particularly natural here: it captures the exploration-exploitation tradeoff agents face when choosing how to forecast, balancing simplicity, reliance on strategies that have worked against openness to alternatives when conditions change.

Our model offers an alternative to explanations based on structural change or nonlinearity. \citet{benigno2023baaack} argue that the post-COVID inflation surge reflects a nonlinear Phillips curve that becomes steep at high output gaps, since inflation remained low during the Great Moderation largely because the economy rarely overheated. Our framework generates similar predictions but through a different mechanism. Specifically, the Phillips curve appears nonlinear because expectation anchoring varies endogenously with the inflation environment, not because the structural relationship itself is nonlinear. This distinction matters for policy because if the curve is structurally nonlinear, the cost of allowing large output gaps is disproportionately high; if anchoring is endogenous, policymakers must also manage credibility.

Finally, our approach differs from standard adaptive learning models \citep{evans2001learning, orphanides2005imperfect} in that agents switch between qualitatively different forecasting rules rather than updating parameters within a single model. Unlike Evans-Honkapohja learning, where coefficients update continuously via recursive least squares, our framework features discrete rule-switching. This produces sharper regime transitions, since beliefs drift gradually in adaptive learning, whereas in our model agents abruptly change strategies when relative performance crosses a threshold. It also introduces asymmetry because switching away from a failing strategy is fast, but rebuilding trust in the anchor requires sustained evidence. Empirically, the models make different predictions about transition speed: adaptive learning predicts smooth adjustment over 10--20 years, while our model predicts faster switches observable in SPF dispersion dynamics during de-anchoring episodes. Relative to \citet{gabaix2020}, whose cognitive discount factor is fixed, our credibility parameter $\theta_t$ is endogenous. Specifically, in Gabaix, the ``attention wedge'' should be constant across monetary regimes; in our model, the effective discount varies with policy credibility, as agents pay more attention to the target when it has performed well. 

Our work also relates to \citet{branch2009}, who study heterogeneous expectations in New Keynesian models. In their framework, a fixed fraction of agents are rational while the remainder use an adaptive rule that extrapolates from past inflation ($\mathbb{E}[\pi] = \theta\pi_{t-1}$), with the parameter $\theta$ controlling the degree of extrapolation. The key difference is that their non-rational agents always use the same backward-looking rule, whereas ours switch between qualitatively different forecasting strategies (anchoring to the central bank's target or extrapolating from recent inflation) depending on which has performed better. This rule-switching is what generates the endogenous credibility dynamics central to our results. Crucially, this switching also generates time-varying inflation persistence, as shocks dissipate quickly when expectations are anchored to the target but propagate through backward-looking dynamics when credibility erodes. 

A feature of our bounded rationality agents is that they perform no reasoning about the underlying structure of the economy. They do not solve optimization problems, form beliefs about other agents, or attempt to infer the central bank's reaction function. Instead, they observe publicly available data (realized inflation and the announced target) and apply simple heuristics. The backward-looking rule, in particular, should not be interpreted as mechanical indexation. Rather, it represents a cognitively cheap heuristic: ``adjust expectations by recent inflation.'' This rule avoids the substantial information processing required for full optimization (forming expectations about demand, competitors, input costs, and monetary policy). When inflation is stable, backward-looking forecasting approximates optimal behavior at minimal cognitive cost; the rule only fails significantly when inflation dynamics shift, precisely when agents should be paying attention.

The specific rules we consider (anchoring to the target or extrapolating from recent inflation) represent one implementation of this modeling approach. Section 4 shows that the framework naturally accommodates additional heuristics, such as trend-following, to address more extreme phenomena like hyperinflation. The intuition is a ``gear-shift'' metaphor. Agents try to forecast inflation as simply as possible, but when simple models start costing them, they shift toward progressively more complex heuristics, from a constant target, to backward-looking behavior, to trend extrapolation, and so on. These cognitive shifts feed back into inflation dynamics. As agents adopt more complex forecasting models, inflation itself becomes more complex, generating phenomena such as endogenous persistence and self-amplifying dynamics, where expectations of rising inflation validate themselves.

The remainder of this paper is organized as follows. Section 2 develops the adaptive-hybrid New Keynesian model, presenting the learning mechanism, the role of key parameters, and supporting evidence on time-varying inflation persistence. Section 3 applies the framework to historical episodes, including the oil shocks, the Great Moderation, the missing disinflation, and the post-COVID surge, and derives policy implications, notably state-dependent optimal policy and the role of credibility as a buffer. Section 4 extends the framework to include trend-following expectations and experience-weighted memory, showing how the same modeling approach addresses more extreme phenomena such as hyperinflation, deflation traps, and stabilization dynamics. Section 5 concludes. Appendices provide details on the multi-armed bandit algorithm, a rational inattention foundation for the simplicity bias, numerical solution methods, robustness analysis, and an alternative Bayesian microfoundation.

%%%%%%%%%%%%%%%%%%%%%%%%%%%%%%%%%%%%%%%%%%%%%%%%%%%%%%%%%%%%%%%%%%%%%%%%%%%%%%%%%%%%%%%%%%%%%%%%%%%%%%%%%%%%%%%%%%%%%%%%%%%%%%%%%%%%%%%%%%%%%%%%%%%%%%%
%%%%%%%%%%%%%%%%%%%%%%%%%%%%%%%%%%%%%%%%%%%%%%%%%%%%%%%%%%%%%%%%%%%%%%%%%%%%%%%%%%%%%%%%%%%%%%%%%%%%%%%%%%%%%%%%%%%%%%%%%%%%%%%%%%%%%%%%%%%%%%%%%%%%%%%
%%%%%%%%%%%%%%%%%%%%%%%%%%%%%%%%%%%%%%%%%%%%%%%%%%%%%%%%%%%%%%%%%%%%%%%%%%%%%%%%%%%%%%%%%%%%%%%%%%%%%%%%%%%%%%%%%%%%%%%%%%%%%%%%%%%%%%%%%%%%%%%%%%%%%%%
%%%%%%%%%%%%%%%%%%%%%%%%%%%%%%%%%%%%%%%%%%%%%%%%%%%%%%%%%%%%%%%%%%%%%%%%%%%%%%%%%%%%%%%%%%%%%%%%%%%%%%%%%%%%%%%%%%%%%%%%%%%%%%%%%%%%%%%%%%%%%%%%%%%%%%%
%%%%%%%%%%%%%%%%%%%%%%%%%%%%%%%%%%%%%%%%%%%%%%%%%%%%%%%%%%%%%%%%%%%%%%%%%%%%%%%%%%%%%%%%%%%%%%%%%%%%%%%%%%%%%%%%%%%%%%%%%%%%%%%%%%%%%%%%%%%%%%%%%%%%%%%

\section{The Model}
\label{sec:model}

The New Keynesian Phillips Curve relates current inflation to expected future inflation, the output gap, and cost-push disturbances. Since expectations enter the inflation equation directly, how agents form them fundamentally shapes inflation dynamics. The Full-Information Rational Expectations (FIRE) benchmark assumes agents know the true model and form expectations consistent with it. This approach is analytically convenient but empirically problematic, since it implies that inflation expectations should jump immediately to the new equilibrium path following any shock, generating sharp but short-lived responses. In practice, inflation exhibits substantial persistence that pure forward-looking models struggle to match \citep{fuhrer1995}.

To address this, \citet{gali1999} developed a hybrid specification incorporating backward-looking elements, finding both forward-looking and backward-looking components to be statistically significant. Large-scale DSGE models such as \citet{smets2007} subsequently adopted similar indexation mechanisms to match the persistence observed in the data. These hybrid specifications improve empirical fit, but they treat the degree of backward-looking behavior as a fixed parameter. The data suggest otherwise: inflation persistence itself varies over time, appearing high during the 1970s,low during the Great Moderation \citep{beechey2012}, and elevated again after 2021 \citep{coibion2025inflation} - the type of variation that models with constant backward-looking weights cannot capture.

Survey evidence suggests that expectation formation itself varies with the inflation environment. \citet{weber2025} document that agents are largely inattentive to inflation when it is low and stable, but become highly attentive when inflation rises. Crucially, this heightened attention does not produce convergence: \citet{mankiw2003disagreement} and \citet{coibion2015} find that forecast disagreement rises during inflationary episodes, suggesting that agents, once attentive, gravitate toward different heuristics rather than toward a common rational forecast. Moreover, \citet{coibion2025inflation} argue that this attention often reveals policy shortcomings, further unanchoring expectations, a self-reinforcing dynamic where attention itself accelerates de-anchoring. This motivates a framework where the degree of anchoring is itself endogenous, earned through forecasting performance rather than assumed.

We develop a model where agents choose between two simple forecasting rules: an anchored rule that trusts the central bank's inflation target, and a backward-looking rule that extrapolates from recent experience. Agents evaluate these rules based on their forecast accuracy and gradually shift toward whichever has performed better. The fraction of agents using the anchored rule, denoted $\theta_t$, becomes a measure of central bank credibility. When the central bank successfully stabilizes inflation, anchored forecasts outperform, credibility rises, and expectations remain anchored. When inflation deviates persistently, backward-looking forecasts gain credibility, anchoring erodes, and expectations become self-fulfilling. This feedback loop can account for the state-dependent dynamics that models with fixed backward-looking weights struggle to capture.

\paragraph{Choosing Between Forecasting Rules}

To model how agents choose between forecasting rules, we draw on reinforcement learning, which studies how agents learn to select from options through trial and experience. The motivation follows from the evidence above: the inflation environment itself shapes expectation formation, with agents reconsidering their approach when inflation becomes volatile. Among the frameworks offered by this literature, we adopt the Multi-Armed Bandit (MAB) framework \citep{sutton2018reinforcement} for its simplicity and intuitive fit.

The MAB framework takes its name from a gambler facing several slot machines, each with an unknown probability of paying out. The gambler must decide which machine to play based on past winnings: exploit the machine that has paid best so far, or explore others that might be better. Crucially, the gambler does not need to understand why one machine pays better than another. She simply tracks which one has worked and gravitates toward it. This captures how bounded rationality agents might approach forecasting: rather than solving a complex optimization problem or forming beliefs about the economy's structure, they observe which forecasting rule has been more accurate and shift toward it.

In our setting, the two 'slot machines' are forecasting rules: agents can anchor expectations to the central bank's target ($\mathbb{E}[\pi_{t+1}] = \pi^*$) or extrapolate from last period's inflation ($\mathbb{E}[\pi_{t+1}] = \pi_{t-1}$) \footnote{The timing within each period is as follows. At the start of period $t$, agents form expectations using $\theta_{t-1}$ and $\pi_{t-1}$, both determined in the previous period. Inflation $\pi_t$ then realizes. Agents observe $\pi_t$, evaluate forecast losses, and update to $\theta_t$, which governs expectations in period $t+1$. The backward-looking rule therefore always uses the most recent inflation available at the time expectations are formed.}. Neither rule requires understanding the economy's structure, as agents need only observe publicly available data (the announced target and realized inflation) and compare forecast errors. Each rule's performance is measured by its squared forecast error:

\begin{align}
L^{CB}_t &= (\pi_t - \pi^*)^2\\
L^{BL}_t &= (\pi_t - \pi_{t-1})^2
\end{align}
  
The choice between these rules has a natural economic interpretation. When inflation is under control, the central bank's target serves as a good proxy for the inflation process. Agents who forecast using the target are, in effect, trusting the central bank to succeed in its mandate. Their forecasts are accurate precisely because monetary policy is working: inflation stays near target, and anchored expectations are validated.

When inflation rises persistently, whether due to a long-lived supply shock or sustained accommodative monetary policy, the target becomes a poor guide. Agents who remain anchored accumulate systematic forecast errors, while the backward-looking rule tracks realized inflation more closely. Observing this performance gap, agents rationally switch their forecasting approach. De-anchoring, in this sense, is not irrational: it reflects agents updating their beliefs about which information source is reliable, based on the central bank's track record of delivering on its stated objective.

So far we have interpreted $\theta_t$ as the fraction of agents using the CB-anchored rule. An alternative interpretation, common in the reinforcement learning literature, views $\theta_t$ not as a population share but as a representative agent's confidence weight on two noisy signals: the central bank's target and recent inflation. In this view, agents form expectations as:

\begin{equation}
\mathbb{E}_t[\pi_{t+1}] = \theta_{t-1} \cdot \pi^* + (1-\theta_{t-1}) \cdot \pi_{t-1}
\end{equation}

The updating rule describes how confidence shifts toward whichever signal has proven more reliable, a form of gradient learning. This interpretation is mathematically equivalent to the population-switching story; we adopt the latter for concreteness, but explore the signal-weighting approach further in Appendix \ref{app:bayesian}.

\paragraph{The Dynamics of Rule Switching}

A common feature in MAB models is to allow agents to update their beliefs only partially each period rather than jumping immediately to the best-performing option. One period of outperformance could reflect noise rather than true superiority, so gradual adjustment guards against overreaction. This is governed by a learning rate parameter $\eta \in (0,1)$. When $\eta$ is small, beliefs adjust slowly; when $\eta$ is large, agents respond quickly to new information. Letting $\theta_t$ denote the current weight on the CB-anchored rule and $\theta^*_t \in 0,1$ the rule that performed best in period $t$, the basic law of motion is:
  
\begin{equation}
\theta_{t+1} = (1-\eta)\theta_t + \eta \cdot \theta^*_t
\end{equation}

In our population share interpretation, $\eta$ measures the fraction of agents reconsidering their forecasting approach each period. This echoes a common modeling device in the New Keynesian literature, specifically the Calvo adjustment mechanism. Just as Calvo pricing assumes a fraction of firms update prices each period while the rest maintain previous prices, our framework assumes a fraction of agents reconsider their forecasting rule while the rest continue with their current approach. This partial adjustment is also consistent with a substantial literature documenting sticky information in expectations. \citet{mankiw2002sticky} show that information diffuses slowly through the population, a finding reinforced by \citet{coibion2015information}, who document significant sluggishness in expectation adjustment.

Substituting the loss functions defined earlier, the law of motion becomes:

\begin{equation}
\theta_{t+1} =
\begin{cases}
(1-\eta)\theta_t + \eta & \text{if } L^{CB}_t < L^{BL}_t  \\
(1-\eta)\theta_t & \text{if } L^{CB}_t > L^{BL}_t
\end{cases}
\end{equation}

When the CB-anchored rule outperforms, all reconsidering agents adopt it, pushing $\theta$ toward one. When the backward-looking rule wins, reconsidering agents switch away from the anchor, pushing $\theta$ toward zero. We explore the edge-case of a tie between the two rules further in the section.

Figure~\ref{fig:learning_mechanism} illustrates this mechanism using a stylized simulation. The left axis tracks inflation (solid black line) as it fluctuates around the central bank's 2\% target, marked by the dashed horizontal line. The gray band indicates a ±1 percentage point range around the target. The right axis tracks credibility $\theta$ (dashed blue line), showing how the fraction of agents using the CB-anchored rule evolves over time. Background shading indicates which forecasting rule is winning in each period: green when inflation is inside the band (CB-anchored rule forecasts accurately), red when inflation is outside (backward-looking rule dominates).

\begin{figure}[htbp]
\centering
\includegraphics[width=\textwidth]{../output/figures/section_2/figure_01_learning_mechanism.png}
\caption{The learning mechanism. Inflation (solid line, left axis) fluctuates around the 2\% target. When inflation remains within the band (light shading), the CB-anchored rule forecasts more accurately and credibility $\theta$ (dashed line, right axis) rises. When inflation deviates persistently (dark shading), the backward-looking rule outperforms and credibility erodes. The gradual adjustment of $\theta$ reflects the learning rate $\eta = 0.10$.}
\label{fig:learning_mechanism}
\end{figure}

The figure reveals three phases. In the first phase (quarters 0--15), inflation remains close to target. The CB-anchored rule forecasts accurately, the background is predominantly light, and credibility stays high near $\theta = 0.95$. In the second phase (quarters 15--35), a persistent shock pushes inflation above the band. The backward-looking rule now tracks inflation more closely than the fixed target, so it wins the performance comparison. The background turns dark, and credibility erodes gradually as a fraction $\eta$ of agents switch away from the anchor each period. By quarter 35, credibility has fallen to approximately $\theta = 0.4$. In the third phase (quarters 35--60), inflation returns to target. The CB-anchored rule regains its accuracy advantage, the background turns light again, and credibility begins to rebuild.

Two features deserve emphasis. First, $\theta$ responds with a lag to changes in the inflation environment. Credibility does not collapse immediately when inflation leaves the band; it erodes gradually as agents accumulate evidence that the backward-looking rule is outperforming. This sluggishness, governed by $\eta$, matches the gradual adjustment of survey expectations documented by \citet{coibion2015information}. Second, the figure hints at an asymmetry that will become central: credibility falls from 0.95 to 0.4 in roughly 20 quarters, but even after 25 quarters of stability it has only recovered to about 0.7. De-anchoring is faster than re-anchoring. We return to this asymmetry after introducing the remaining model elements.

\paragraph{A Simplicity Bias}

In Figure~\ref{fig:learning_mechanism}, the CB-anchored rule wins whenever inflation falls inside the band, even when the backward-looking forecast is slightly more accurate. This reflects an important feature of the model: a tie-breaking bias toward simpler forecasting rules.

The CB-anchored rule requires recalling a single fixed number: the announced target. The backward-looking rule requires tracking a time-varying process, continuously monitoring and updating based on realized inflation. Following the rational inattention literature \citep{sims2003implications}, we assume this additional cognitive effort carries a cost. Agents will only incur it when the accuracy gain is sufficiently large. This is consistent with \citet{mackowiak2009optimal}, who show that agents with finite processing capacity optimally allocate attention toward the most variable signals: when inflation is stable, it does not warrant the cognitive cost of active tracking, and the CB-anchored rule dominates; when inflation becomes volatile, the payoff to attention rises and agents rationally switch to the more demanding backward-looking rule. This yields a modified switching rule:
  
\begin{equation}
\label{eq:switching_rule}
\theta_{t+1} =
\begin{cases}
(1-\eta)\theta_t + \eta & \text{if } L^{CB}_t \leq L^{BL}_t + \varepsilon\\
(1-\eta)\theta_t & \text{if } L^{CB}_t > L^{BL}_t + \varepsilon
\end{cases}
\end{equation}

where $\varepsilon > 0$ represents the threshold below which agents prefer the simpler rule. The CB-anchored rule now wins not only when it forecasts better, but also when it forecasts nearly as well. Appendix~\ref{app:entropy} grounds this parameter in the rational inattention literature, showing that $\varepsilon = 10^{-4}$, equivalent to forgiving a forecast that misses by one percentage point, is consistent with estimated information costs.

This parameter is not merely a technical convenience; it is essential for the model's dynamics. Consider what happens without it ($\varepsilon = 0$). Suppose credibility has eroded during an inflationary episode, and the central bank successfully returns inflation to target gradually. At this point, with $\pi_t = \pi^*$, both rules generate identical forecasts: the CB-anchored rule forecasts $\pi^*$ by construction, while the backward-looking rule forecasts $\pi_{t-1} = \pi^*$. With identical forecasts come identical losses. Agents have no reason to switch back to the anchor. Credibility remains stuck near zero despite the central bank's success. This is the re-anchoring problem: without $\varepsilon > 0$, credibility once lost cannot be regained through good policy alone. Appendix~\ref{app:entropy} illustrates this with simulations.

The simplicity bias resolves this problem. When inflation stabilizes at target, both rules are equally accurate, but the CB-anchored rule is cognitively cheaper. The tie-breaker favors it, and credibility gradually rebuilds. In this sense, $\varepsilon$ captures the intuition that agents will default to simpler heuristics when more complex approaches offer no clear advantage.

\paragraph{Memory and Cumulative Performance}

Credibility is typically thought to be built gradually and lost only through sustained policy failure \citep{blinder2000}. Yet in Figure~\ref{fig:learning_mechanism}, a single quarter of elevated inflation is enough to tip the comparison in favor of the backward-looking rule. A single bad outcome should not undo years of successful stabilization.

To address this, we allow agents to evaluate forecast performance over a memory window of $k$ periods rather than a single quarter. The loss functions become cumulative:
  
\begin{align}
L^{CB}_t &= \sum_{s=t-k+1}^{t} (\pi_s - \pi^*)^2 \\
L^{BL}_t &= \sum_{s=t-k+1}^{t} (\pi_s - \pi_{s-1})^2
\end{align}
  
Agents now compare track records over the recent past, not just the most recent observation. The switching rule~\eqref{eq:switching_rule} remains unchanged, with $L^{CB}_t$ and $L^{BL}_t$ now representing cumulative losses over the $k$-period window.

Figure~\ref{fig:k_epsilon_roles} illustrates the roles of both $\varepsilon$ and $k$. Panel (a) shows the re-anchoring problem: without the simplicity bias ($\varepsilon = 0$), a soft landing to target leaves both rules equally accurate, and credibility remains stuck near zero. With $\varepsilon > 0$, the tie-breaker favors the CB-anchored rule, and credibility rebuilds. Panel (b) shows how memory length affects resilience to shocks. With $k = 1$, even a brief single-quarter spike outside the band triggers de-anchoring. With $k = 3$, the same spike is averaged with quarters of stability, the CB-anchored rule retains its cumulative advantage, and credibility is preserved. Only when the deviation persists for several quarters does the backward-looking rule establish a superior track record.

\begin{figure}[htbp]
\centering
\includegraphics[width=\textwidth]{../output/figures/section_2/figure_02_k_epsilon_roles.png}
\caption{The roles of $\varepsilon$ and $k$. Panel (a): Without the simplicity bias ($\varepsilon = 0$), a soft landing to target leaves credibility stuck near zero. During the gradual return to target, inflation is always closer to its previous value than to the target, so the backward-looking rule keeps winning (dark shading) and credibility cannot rebuild. Panel (b): With $k = 1$ (dark blue, dashed), a single-quarter spike triggers immediate de-anchoring. With $k = 3$ (light blue, solid), the brief deviation is outvoted by surrounding stable periods, and only the persistent 12-quarter departure erodes credibility.}
\label{fig:k_epsilon_roles}
\end{figure}

The memory window has a symmetric consequence: it makes credibility both harder to lose and harder to regain. When inflation rises, the CB-anchored rule must underperform for several consecutive periods before agents switch. But when inflation returns to target, the backward-looking rule retains its cumulative advantage from the recent inflationary period. Agents remain de-anchored until enough low-inflation quarters accumulate to tip the comparison back. This creates endogenous persistence: expectations stay backward-looking for an extended period after the shock has passed, sustaining elevated inflation even as the underlying disturbance fades. We explore this mechanism further in Section~\ref{sec:applications}.

\paragraph{Rational Agents in the Population}

The model so far assumes that all agents use the heuristics described above. This is a useful benchmark, but survey evidence suggests a more nuanced picture. Professional forecasters, for instance, produce expectations that are closer to rational benchmarks than household surveys \citep{coibion2015information}. Financial market participants have strong incentives to form accurate forecasts and access to sophisticated models. This suggests that some fraction of the population may form expectations in a manner closer to the Full-Information Rational Expectations (FIRE) benchmark.

We incorporate this by assuming a fraction $\lambda \in [0,1]$ of agents form FIRE expectations, while the remaining $(1-\lambda)$ use the adaptive learning mechanism described above. Aggregate expectations become a weighted average:
  
\begin{equation} 
\label{eq:aggregate_expectations}
\mathbb{E}_t[\pi_{t+1}] = \lambda \mathbb{E}_t^{FIRE}[\pi_{t+1}] + (1-\lambda)\left[\theta_{t-1} \pi^* + (1-\theta_{t-1})\pi_{t-1}\right]
\end{equation}

When $\lambda = 1$, this reduces to the standard New Keynesian model with rational expectations. When $\lambda = 0$, we recover the pure adaptive model. Intermediate values capture an economy where sophisticated agents coexist with those using simpler forecasting rules.

Introducing FIRE agents raises a modeling question. Are FIRE agents "sophisticated," aware of the bounded rationality agents in the population, or are they "naive," solving as if all agents were rational? The naive specification simplifies the problem considerably: FIRE expectations are formed by solving the standard New Keynesian model assuming $\lambda = 1$, then aggregated ex-post with the MAB learners' expectations. The sophisticated specification is significantly more challenging: FIRE agents observe the credibility state $\theta_t$ and solve forward taking its evolution into account. This creates a fixed-point problem, since FIRE forecasts depend on the path of credibility, which depends on realized inflation, which in turn depends on aggregate expectations including the FIRE component.

We adopt the sophisticated specification throughout to isolate the MAB mechanism as the sole source of bounded rationality in the model. The naive case remains useful, however, because it provides a good approximation in most states, we use it as a warm start for the iterative algorithm that solves the sophisticated model. Appendix~\ref{app:computational} describes the solution method in detail.

We close the model with the standard equations of the New Keynesian framework \citep{gali2015monetary}. The Phillips Curve relates current inflation to expected future inflation and the output gap:

\begin{equation}
\pi_t = \beta \mathbb{E}_t[\pi_{t+1}] + \kappa y_t + u_t
\end{equation}

where $\beta$ is the discount factor, $\kappa$ is the slope of the Phillips Curve, and $u_t$ is a cost-push shock, and the expectation term $\mathbb{E}_t[\pi_{t+1}]$ follows equation~\eqref{eq:aggregate_expectations}, combining FIRE and MAB forecasts. The IS curve captures intertemporal consumption smoothing:

\begin{equation}
y_t = \mathbb{E}_t[y_{t+1}] - \sigma(i_t - \mathbb{E}_t[\pi_{t+1}] - r^n)
\end{equation}

where $y_t$ is the output gap, $i_t$ is the nominal interest rate, $\sigma$ is the intertemporal elasticity of substitution (equal to $1/\sigma$ in the Gal\'i notation, where $\sigma$ denotes the coefficient of relative risk aversion; the two coincide at our benchmark of log utility), and $r^n$ is the natural rate of interest. Monetary policy follows a Taylor rule:

\begin{equation}
i_t = r^n + \pi^* + \phi_\pi(\pi_t - \pi^*) + \phi_y y_t
\end{equation}

where $\phi_\pi$ and $\phi_y$ govern the policy response to inflation deviations and the output gap.

One modeling choice deserves emphasis. While bounded rationality agents use heuristics to forecast inflation, we assume all agents form rational expectations over output, including those using the MAB learning mechanism. That is, the expectations in the IS curve are FIRE for all agents. This is an identification restriction: endogenizing both inflation and output expectations simultaneously through the credibility mechanism would confound the source of the results, making it impossible to attribute dynamics cleanly to the inflation credibility channel. By holding output expectations at the rational benchmark, every departure from the standard model traces directly to the evolution of $\theta_t$.

Finally, we set the inflation target to $\pi^* = 2$, matching the Federal Reserve's official target. Throughout the paper, inflation rates and targets are expressed in annualized percentage points; the code operates in quarterly decimal rates (so $\pi^* = 0.005$), and all reported results are converted to annualized percent for presentation. The standard New Keynesian model has a zero-inflation steady state; we add a constant term $(1-\beta)\pi^*$ to the Phillips Curve and include $\pi^*$ in the Taylor rule intercept so that the economy converges to 2\% rather than zero. This allows us to discuss credibility in intuitive terms: anchoring means expecting inflation to return to 2\%.

Figure~\ref{fig:lambda_comparison} illustrates how the fraction of rational agents affects the economy's response to cost-push shocks. The left column shows a single persistent shock (0.5\%, $\rho = 0.8$); the right column shows repeated shocks (0.75\% for 12 quarters). The top row displays inflation paths; the bottom row shows credibility dynamics. Each panel compares three cases: all adaptive agents ($\lambda = 0$), a mixed population ($\lambda = 0.3$), and all rational agents ($\lambda = 1$).

\begin{figure}[htbp]
\centering
\includegraphics[width=\textwidth]{../output/figures/section_2/figure_03_lambda_comparison.png}
\caption{Effect of rational agent fraction ($\lambda$) on inflation and credibility dynamics under cost-push shocks.}
\label{fig:lambda_comparison}
\end{figure}

The two panels reveal different dynamics depending on shock duration. Under a single persistent shock (left), FIRE agents anticipate the full shock path and generate the largest inflation response; the adaptive models, anchored by the simplicity bias, produce lower peaks with credibility intact throughout. Under sustained shocks (right), the pattern reverses: the mixed population ($\lambda = 0.3$) generates the highest peak inflation, exceeding both the pure adaptive and pure rational cases. This occurs because FIRE agents push up aggregate inflation, which causes the adaptive agents to de-anchor, and their backward-looking expectations then amplify and prolong the shock. The pure adaptive economy ($\lambda = 0$) also de-anchors but lacks the FIRE-driven initial amplification; the pure rational economy ($\lambda = 1$) has no credibility mechanism to generate persistence. The hybrid case thus represents the scenario where rational expectations trigger the initial inflation that adaptive learning then propagates.

\paragraph{Calibration} We now examine the model's aggregate dynamics using parameters calibrated from established sources in the New Keynesian and behavioral macroeconomics literatures. Table~\ref{tab:calibration} summarizes the baseline calibration; we discuss the rationale for each parameter below.

\begin{table}[htbp]
\centering
\caption{Baseline Calibration}
\label{tab:calibration}
\begin{tabular}{llcl}
\hline
Parameter & Description & Value & Source \\
\hline
\multicolumn{4}{l}{\textit{Structural Parameters}} \\
$\beta$ & Discount factor & 0.99 & Gal\'i (2015) \\
$\sigma$ & Intertemporal elasticity & 1.0 & Log utility benchmark \\
$\kappa$ & Phillips curve slope & 0.024 & Gal\'i \& Gertler (1999) \\
$\phi_\pi$ & Taylor rule: inflation & 1.5 & Taylor (1993) \\
$\phi_y$ & Taylor rule: output & 0.125 & Taylor (1993) \\
$\pi^*$ & Inflation target & 2\% & Federal Reserve mandate \\
\hline
\multicolumn{4}{l}{\textit{Behavioral Parameters}} \\
$\lambda$ & FIRE fraction & 0.35 & Calibration discussion \\
$\eta$ & Learning speed & 0.10 & Malmendier \& Nagel (2016) \\
$k$ & Memory window & 3 & -- \\
$\varepsilon$ & Tie-breaking threshold & $10^{-4}$ & Appendix \ref{app:entropy} \\
\hline
\end{tabular}
\end{table}

The discount factor $\beta = 0.99$ corresponds to a 4\% annual real interest rate, standard in quarterly New Keynesian models \citep{gali2015monetary}. The intertemporal elasticity of substitution $\sigma = 1$ (log utility) is a common benchmark. The Taylor rule coefficients $\phi_\pi = 1.5$ and $\phi_y = 0.125$ match \citet{taylor1993discretion} and satisfy the Taylor principle. The Phillips curve slope $\kappa = 0.024$ is at the lower end of empirical estimates, consistent with the ``flat Phillips curve'' literature \citep{gali1999, hazell2022slope}. Our choice reflects the view that aggregate Phillips curve estimates are biased by expectation anchoring (precisely the mechanism our model endogenizes).

The FIRE fraction $\lambda = 0.35$ implies roughly one-third of agents form model-consistent expectations. No single paper estimates this parameter directly; instead, the value is informed by converging evidence from several sources. \citet{coibion2015} find that professional forecasters exhibit near-rational behavior while household expectations deviate substantially, suggesting a minority of agents approximate the FIRE benchmark. \citet{mankiw2003disagreement} document persistent disagreement across forecasters inconsistent with common-information rational expectations, pointing to fundamental heterogeneity in forecasting sophistication. \citet{branch2004model} estimate that a non-negligible fraction of agents use sophisticated forecasting models in experimental settings. Taken together, these findings support a FIRE share in the range of one-quarter to one-half; we set $\lambda = 0.35$ as a central value. The learning speed $\eta = 0.10$, implying agents reconsider their forecasting rule roughly every 2.5 years on average, is chosen so that regime transitions unfold over multi-year horizons, consistent with episodes like the Volcker disinflation. \citet{malmendier2016learning} provide direct support, finding that experienced inflation shapes expectations with a half-life of several years; \citet{weber2025} show that attention to inflation is state-dependent, rising sharply during inflationary episodes.

The memory window $k = 3$ quarters reflects evidence that agents evaluate forecast performance over multiple periods rather than reacting to single observations \citep{blinder2000}. The tie-breaking threshold $\varepsilon = 10^{-4}$ is equivalent to tolerating a one-percentage-point forecast error before switching to the more complex rule. Appendix~\ref{app:entropy} shows this value is consistent with information costs estimated in the rational inattention literature \citep{sims2003implications, mackowiak2009optimal}. Results are qualitatively robust to reasonable variations in these parameters.

We simulate the model under three distinct shock regimes with initial credibility calibrated to each era. The first, labeled "Great Moderation," features moderate shocks ($\sigma_u = 0.0006$) representative of the stable macroeconomic environment from the mid-1980s to 2007; we set $\theta_0 = 0.5$, reflecting moderate credibility at the start of the period. The second, labeled "1970s-Style," features large and volatile shocks sufficient to generate inflation peaks near 10; we set $\theta_0 = 0.9$, showing that even high initial credibility erodes under sufficiently large and persistent shocks. The third, labeled "Credibility Cycle," depicts a complete cycle: the economy begins in calm (years 0--5), experiences a decade of turbulence (years 5--15,  shaded), and returns to tranquility (years 15--30); we set $\theta_0 = 1$ to trace how credibility evolves from a fully anchored starting point. Figure~\ref{fig:shock_regimes} displays the resulting paths.

\begin{figure}[htbp]
\centering
\includegraphics[width=\textwidth]{../output/figures/section_2/figure_04_shock_regimes.png}
\caption{Representative simulations under different shock regimes. Top row: inflation dynamics. Bottom row: credibility ($\theta$) dynamics. Left column: Great Moderation with moderate shocks. Middle column: 1970s-style high volatility. Right column: Credibility cycle showing anchoring, de-anchoring during the turbulent period (shaded), and subsequent re-anchoring.}
\label{fig:shock_regimes}
\end{figure}

The Great Moderation regime reveals a subtle but important property. Despite moderate shock volatility, credibility fluctuates considerably as shocks occasionally cause the backward-looking rule to outperform. However, these fluctuations do not destabilize inflation: when inflation remains close to target, both forecasting rules yield similar predictions, so shifts in composition have negligible effects on aggregate expectations. This self-stabilizing property ensures the good equilibrium is robust.

The 1970s-style regime presents a different picture. Despite starting at high credibility ($\theta_0 = 0.9$), large shocks drive inflation far from target, causing the backward-looking rule to generate forecasts that differ substantially from $\pi^*$. When the BL rule outperforms, agents switch away from the CB anchor and $\theta$ falls. Lower credibility increases the weight on backward-looking expectations, which raises inflation persistence, which makes the backward-looking rule more accurate, triggering further de-anchoring. The result is a self-reinforcing cycle: $\theta$ erodes and remains depressed, and inflation exhibits extreme volatility with prolonged deviations from target. As inflation rises, backward-looking forecasts outperform, agents switch away from the anchor, and their revised expectations validate continued high inflation, creating a de-anchoring spiral.

The credibility cycle illustrates the complete dynamics. During the initial five years of calm, credibility remains anchored near $\theta = 1$. When turbulence begins, credibility falls within the decade as large forecast errors favor the backward-looking rule, reaching a trough around $\theta \approx 0.3$ by year 15. The key insight emerges in the recovery phase, where credibility rebuilds gradually, reaching 0.8 within about 3 years and 0.9 within 5 years of the turbulence ending, but full recovery to $\theta \approx 1$ takes roughly 15 years. The asymmetry is notable: a decade of turbulence erodes credibility from 1.0 to 0.3, while rebuilding from 0.9 to 1.0 alone takes nearly a decade more.

\paragraph{Transitory versus Persistent Shocks}

For transitory shocks of moderate magnitude, the adaptive model behaves similarly to FIRE. Table~\ref{tab:transitory_persistent} demonstrates this using a 4.5\% annualized cost-push shock under varying persistence. When $\rho = 0$, both models produce nearly identical peak inflation, and credibility remains fully intact. At moderate persistence ($\rho = 0.4$), the models begin to diverge, as the adaptive model produces \textit{lower} peak inflation than FIRE, because anchored expectations dampen the shock that rational agents fully anticipate. Credibility remains intact throughout, confirming that moderate persistence alone does not trigger de-anchoring. Only at high persistence ($\rho = 0.85$) does the backward-looking rule establish a superior track record, causing credibility to erode ($\theta$ falls to 0.59). Even then, peak inflation in the adaptive model remains well below FIRE, reflecting the stabilizing effect of the agents who remain anchored. This result clarifies when the credibility mechanism matters: transitory disturbances are absorbed without de-anchoring; it is persistence, not magnitude alone, that triggers credibility erosion.\footnote{The adaptive model's peak inflation decreases slightly with persistence despite $\theta$ remaining near one. This reflects a modeling choice: output expectations are rational for all agents, so the anticipated persistent shock generates a large negative output gap that partially offsets inflationary pressure. The effect disappears once credibility erodes and inflation expectations de-anchor.}

\begin{table}[htbp]
\centering
\caption{Transitory vs Persistent Shocks: Only Persistence Triggers De-Anchoring}
\label{tab:transitory_persistent}
\begin{tabular}{ccccc}
\toprule
Persistence ($\rho$) & Peak $\pi$ (FIRE) & Peak $\pi$ (Adaptive) & Difference & Min $\theta$ \\
\midrule
0.0 & 4.33\% & 4.36\% & 0.03\% & 1.00 \\
0.4 & 7.00\% & 4.27\% & $-$2.73\% & 1.00 \\
0.85 & 22.78\% & 3.76\% & $-$19.02\% & 0.59 \\
\bottomrule
\end{tabular}
\begin{tablenotes}
\small
\item \textit{Notes}: Simulations use a 4.5\% annualized cost-push shock with baseline parameters ($\eta = 0.10$, $\phi_\pi = 1.5$, $k = 3$, $\varepsilon = 10^{-4}$). Peak inflation reported as deviation from target. ``Difference'' is Adaptive minus FIRE.
\end{tablenotes}
\end{table}

\paragraph{Faster Learning Can Backfire}

A second key result concerns the role of learning speed. Table~\ref{tab:paradox_learning} reports outcomes for a large persistent shock ($\rho = 0.85$) under three learning speeds. Contrary to the intuition that faster learning should improve outcomes, faster learners ($\eta = 0.30$) suffer substantially worse inflation: peak inflation more than doubles (14.75\% versus 6.68\%), and the economy fails to return to target within the simulation horizon. Faster learners switch more quickly to backward-looking forecasts when those forecasts outperform. But this responsiveness becomes a liability, as their backward-looking expectations generate the elevated inflation that validates those expectations, trapping the economy in a de-anchored state. Slower learners ($\eta = 0.10$) retain enough agents using the CB-anchored rule to pull inflation back toward target within 28 quarters, breaking the feedback loop before it becomes self-sustaining.

\begin{table}[htbp]
\centering
\caption{The Paradox of Faster Learning: Higher $\eta$ Leads to Worse Outcomes}
\label{tab:paradox_learning}
\begin{tabular}{ccccc}
\toprule
Learning Speed ($\eta$) & Peak $\Delta\pi$ & Min $\theta$ & Quarters to Target & Final $\theta$ \\
\midrule
0.10 & 6.68\% & 0.122 & 28 & 1.000 \\
0.20 & 8.35\% & 0.000 & 96 & 0.995 \\
0.30 & 14.75\% & 0.000 & $>$115 & 0.000 \\
\bottomrule
\end{tabular}
\begin{tablenotes}
\small
\item \textit{Notes}: Simulations use an 8\% annualized persistent shock ($\rho = 0.85$) with baseline parameters ($\phi_\pi = 1.5$, $k = 3$, $\varepsilon = 10^{-4}$). ``Quarters to Target'' counts periods after shock until inflation returns within 0.2\% of target.
\end{tablenotes}
\end{table}

The sharp nonlinearity between $\eta = 0.20$ and $\eta = 0.30$---where return-to-target time jumps from 96 quarters to beyond the simulation horizon---suggests a critical threshold of $\eta$ beyond which the de-anchoring spiral becomes self-sustaining. Characterizing this threshold analytically, and how it depends on the policy stance and shock persistence, is left to future work.

\paragraph{Supporting Evidence: Time-Varying Inflation Persistence} The model predicts that inflation persistence should vary with credibility, rising when expectations are de-anchored and falling when anchored. This pattern is well-documented in the literature \citep{cogley2005drifts, stock2007has}, but we briefly present the evidence here to show it aligns with our framework.

Using rolling-window AR(1) regressions on quarterly U.S. CPI inflation from 1960Q1 through 2024Q4, we estimate persistence $\rho$ for each quarter on a trailing 40-quarter window. Chow tests indicate significant structural breaks at the Volcker appointment (1979Q4), the Great Moderation onset (1984Q1), and the financial crisis (2008Q4). Table \ref{tab:persistence_regimes} summarizes persistence by regime.

\begin{table}[htbp]
\centering
\caption{Inflation Persistence by Monetary Policy Regime}
\label{tab:persistence_regimes}
\begin{tabular}{lcccccc}
\toprule
Regime & Mean $\hat{\rho}$ & Std. Dev. & $N$ & Break Date & $F-stat$ & $p-value$ \\
\midrule
Great Inflation (1968--1979) & 0.729 & 0.042 & 22 & 1979Q4 & 13.13 & $<$0.001 \\
Volcker Era (1980--1984) & 0.671 & 0.080 & 20 & 1984Q1 & 23.77 & $<$0.001 \\
Great Moderation (1985--2007) & 0.255 & 0.319 & 92 & 2008Q4 & 33.55 & $<$0.001 \\
Post-Crisis (2008--2019) & $-0.045$ & 0.105 & 48 & 2020Q1 & 1.93 & 0.147 \\
COVID Era (2020--2024) & 0.392 & 0.221 & 20 & -- & -- & -- \\
\bottomrule
\end{tabular}
\end{table}

\begin{figure}[htbp]
\centering
\includegraphics[width=\textwidth]{../output/figures/section_2/figure_05_persistence.png}
\caption{Time-varying inflation persistence, 1970--2024. Solid line shows rolling 40-quarter AR(1) coefficients; shaded region indicates 95\% confidence interval. Vertical lines mark Chow test candidate dates at major policy events (dashed lines indicate $p < 0.05$; dotted line indicates $p > 0.05$).}
\label{fig:persistence}
\end{figure}

Figure \ref{fig:persistence} shows the results. Persistence averages $\hat{\rho} = 0.73$ during the Great Inflation, remains elevated through the Volcker era ($\hat{\rho} = 0.67$) as credibility had not yet been established, then declines substantially during the Great Moderation ($\hat{\rho} = 0.26$) and falls to essentially zero post-crisis ($\hat{\rho} = -0.05$). The COVID era reverses this trend ($\hat{\rho} = 0.39$), though persistence remains well below 1970s levels. Appendix~\ref{app:robustness} provides a robustness check using a state-space model with Kalman filtering; the two approaches yield highly correlated persistence paths ($\rho = 0.86$), suggesting the results are robust to the estimation method.

This pattern is consistent with what our model predicts: when credibility is high, expectations anchor to target and shocks are absorbed quickly; when credibility erodes, backward-looking behavior dominates and inflation inherits persistence from its own past. The following section uses simulations to explore how these dynamics play out in specific historical episodes.


%%%%%%%%%%%%%%%%%%%%%%%%%%%%%%%%%%%%%%%%%%%%%%%%%%%%%%%%%%%%%%%%%%%%%%%%%%%%%%%%%%%%%%%%%%%%%%%%%%%%%%%%%%%%%%%%%%%%%%%%%%%%%%%%%%%%%%%%%%%%%%%%%%%%%%%
%%%%%%%%%%%%%%%%%%%%%%%%%%%%%%%%%%%%%%%%%%%%%%%%%%%%%%%%%%%%%%%%%%%%%%%%%%%%%%%%%%%%%%%%%%%%%%%%%%%%%%%%%%%%%%%%%%%%%%%%%%%%%%%%%%%%%%%%%%%%%%%%%%%%%%%
%%%%%%%%%%%%%%%%%%%%%%%%%%%%%%%%%%%%%%%%%%%%%%%%%%%%%%%%%%%%%%%%%%%%%%%%%%%%%%%%%%%%%%%%%%%%%%%%%%%%%%%%%%%%%%%%%%%%%%%%%%%%%%%%%%%%%%%%%%%%%%%%%%%%%%%
%%%%%%%%%%%%%%%%%%%%%%%%%%%%%%%%%%%%%%%%%%%%%%%%%%%%%%%%%%%%%%%%%%%%%%%%%%%%%%%%%%%%%%%%%%%%%%%%%%%%%%%%%%%%%%%%%%%%%%%%%%%%%%%%%%%%%%%%%%%%%%%%%%%%%%%
%%%%%%%%%%%%%%%%%%%%%%%%%%%%%%%%%%%%%%%%%%%%%%%%%%%%%%%%%%%%%%%%%%%%%%%%%%%%%%%%%%%%%%%%%%%%%%%%%%%%%%%%%%%%%%%%%%%%%%%%%%%%%%%%%%%%%%%%%%%%%%%%%%%%%%%

\section{Applications and Policy Implications}
\label{sec:applications}

This section applies the adaptive-hybrid framework to episodes and puzzles that are difficult to reconcile within a standard New Keynesian model with fixed expectation parameters, and draws out implications for monetary policy design. The simulations that follow are mechanism illustrations and comparative statics, not moment-matching exercises; their purpose is to show that the model's qualitative predictions---regime-dependent persistence, credibility hysteresis, asymmetric stabilization costs---align with the broad features of each episode. We first examine inflation regimes in U.S. history, such as the 1970s oil shocks, the Great Moderation, the missing disinflation of 2008–09, and the post-pandemic surge, showing how the evolution of $\theta_t$ provides a common mechanism for episodes that are often explained through structural breaks or parameter instability. We then address two empirical puzzles of the Phillips curve: Friedman's "long and variable lags" and the apparent time-variation in reduced-form slope estimates documented by \citet{mavroeidis2014empirical}. Finally, we turn to policy design, exploring how optimal aggressiveness depends on the credibility state, why the case for "looking through" supply shocks weakens as expectations de-anchor, and how the speed of disinflation interacts with credibility recovery.

\paragraph{Inflation Regimes in US History} US inflation since 1970 offers a useful setting to examine our framework, as it spans several distinct regimes that are difficult to reconcile with standard models. The 1970s featured persistent high inflation amplified by oil shocks; the Volcker disinflation (1979--1982) required extreme policy tightening and produced severe output costs; the Great Moderation (1984--2007) saw remarkably stable inflation with muted responses to shocks; the Great Recession (2008--2012) produced the ``missing disinflation'': inflation remained stable despite the largest output collapse since the Depression; and the post-pandemic period (2021--2024) saw inflation that was initially muted despite adverse conditions, then spiked and remained persistently high, before declining without substantial employment costs.

These regime shifts are well-documented. \citet{cogley2005drifts} show inflation persistence shifted from near-unit-root behavior in the 1970s to stationary dynamics thereafter, \citet{stock2007has} document a structural break around 1984, and \citet{hazell2022slope} argue that the apparent variation in Phillips curve slope across eras reflects changes in expectation anchoring rather than in structural parameters. The hybrid Phillips curve of \citet{gali1999} captures backward-looking behavior but treats the backward-looking weight as a fixed parameter, leaving unexplained why persistence was high in the 1970s, low in the 1990s, and elevated again after 2021. Our framework speaks to this pattern, as the evolution of $\theta_t$ provides a mechanism for these regime shifts without requiring changes to structural parameters.

\paragraph{Oil Shocks and the Role of Initial Credibility} Why did oil price shocks cause stagflation in the 1970s but not in 2008, despite similar-sized commodity price increases? \citet{blanchard2007macroeconomic} document this puzzle systematically, showing that the oil shocks of 2007--2008 were comparable in magnitude to those of 1973--1974 and 1979--1980, yet their macroeconomic effects were substantially smaller. They consider several explanations, such as a smaller share of oil in production, more flexible labor markets, and improvements in monetary policy, and conclude that improved central bank credibility played a key role. But what is credibility, and how does it shape inflation dynamics? Our framework offers one interpretation.

In this view, the mechanism operates through expectation anchoring. When credibility is high, most agents forecast inflation near the central bank's target; a cost-push shock raises inflation temporarily, but expectations remain anchored, limiting second-round effects. When credibility is low, most agents extrapolate recent inflation; the same shock rewards their backward-looking forecasts, triggering further de-anchoring and amplifying the initial impulse. This suggests that identical shocks can produce substantially different outcomes depending on the stock of credibility accumulated before the shock hits.

\begin{figure}[htbp]
\centering
\includegraphics[width=0.95\textwidth]{../output/figures/section_3/figure_06_oil_shocks.png}
\caption{Shock transmission under low versus high initial credibility. Both panels show responses to an identical 1 percentage point (quarterly) cost-push shock. Left: inflation peaks at 8.2\% under low credibility ($\theta_0 = 0.2$) versus 5.5\% under high credibility ($\theta_0 = 1.0$). Right: credibility collapses to 0.11 in the low-$\theta_0$ scenario, while remaining at 1.0 in the high-$\theta_0$ case; the rapid return to target is consistent with anchored expectations.}
\label{fig:oil_shocks}
\end{figure}

Figure~\ref{fig:oil_shocks} illustrates this mechanism. We simulate an identical 1 percentage point (quarterly) cost-push shock under two initial conditions: $\theta_0 = 0.2$ (low credibility, representing the 1970s) and $\theta_0 = 1.0$ (high credibility, representing 2008). The monetary policy rule is held constant at $\phi_\pi = 1.5$ in both scenarios; the only difference is initial credibility. Under low credibility, inflation peaks at 8.2\% and remains above 3\% for 12 quarters; credibility collapses to 0.11 as the shock rewards backward-looking forecasts. Under high credibility, the same shock produces a peak of only 5.5\%, inflation returns below 3\% within 4 quarters, and credibility remains at 1.0 throughout; the swift return to target rewards the CB-anchored forecasting rule, making high credibility self-reinforcing. The 2.7 percentage point difference in peak inflation and the threefold difference in persistence arise entirely from the state of expectation anchoring when the shock hits. 

\paragraph{Three Episodes: Great Moderation, Missing Disinflation, and Post-Pandemic Inflation} We now apply this framework to three episodes that span the range of credibility regimes. For each, we simulate the model with baseline parameters ($\lambda = 0.35$, $\eta = 0.10$, $\kappa = 0.024$) and compare against the standard New Keynesian benchmark to isolate the role of endogenous credibility dynamics.

The Great Moderation (1984–2007) has been attributed to improved monetary policy, structural changes, and smaller shocks \citep{stock2003has, bernanke2004great}. The post-COVID inflation surge casts doubt on the first two explanations: large shocks returned, and so did inflation, suggesting the earlier stability was not purely structural. Our framework suggests a complementary mechanism: high credibility inherited from the calm period following the Volcker disinflation acted as an automatic stabilizer, anchoring expectations and preventing the de-anchoring spirals that characterized the 1970s. We simulate this regime with high initial credibility ($\theta_0 = 0.95$) facing small, low-persistence shocks (cost-push shocks drawn from a normal distribution with standard deviation of 0.2 percentage points and persistence $\rho = 0.3$). The prediction is that when credibility is high, small shocks should not trigger de-anchoring, and the adaptive model should track the standard NK  benchmark closely, with muted responses, since anchored agents continue forecasting near target regardless of realized shocks.

\begin{figure}[htpb!]
\centering
\includegraphics[width=\textwidth]{../output/figures/section_3/figure_07_great_moderation.png}
\caption{The Great Moderation (1984--2007). Inflation paths for the adaptive model and standard NK track closely (correlation = 0.94), as high credibility keeps expectations anchored throughout.}
\label{fig:great_moderation}
\end{figure}

Figure \ref{fig:great_moderation} shows that inflation paths under the adaptive and standard NK models are nearly indistinguishable, with a correlation of 0.94. Credibility remains near 1.0 throughout the simulation, and the difference between models rarely exceeds 0.2 percentage points. With high $\theta$, the adaptive model's aggregate expectation is dominated by the anchored component, bringing aggregate dynamics close to the standard NK benchmark. The Great Moderation, in this interpretation, was a period where credibility was high enough that the distinction between our model and standard NK became empirically negligible. This interpretation complements rather than replaces the ``good policy, good luck, or structural change'' debate \citep{stock2003has}. Even if smaller shocks and improved policy were primary drivers, high credibility amplified their stabilizing effects by ensuring that the shocks which did occur were absorbed by expectations rather than propagated through backward-looking dynamics. The model suggests a virtuous cycle: successful stabilization builds credibility, which makes future stabilization easier, which builds more credibility.

Despite the largest output collapse since the Great Depression, inflation during the Great Recession declined substantially less than standard models predicted, a pattern known as the "missing disinflation" puzzle \citep{ball2011, coibion2015}. Our framework suggests that high credibility accumulated during the Great Moderation insulated expectations from the demand collapse. To illustrate, we simulate a negative 0.5 percentage point (quarterly) cost-push shock hitting an economy with high inherited credibility ($\theta_0 = 0.95$) and persistence $\rho = 0.7$. Because $\theta$ remains elevated, aggregate expectations stay anchored near target even as realized inflation falls, dampening the deflationary feedback loop that standard models would predict.

\begin{figure}[htpb!]
\centering
\includegraphics[width=\textwidth]{../output/figures/section_3/figure_08_missing_disinflation.png}
\caption{The Missing Disinflation (2008--2012). Standard NK predicts deflation ($-3.7\%$), while the adaptive model with high credibility shows a much smaller decline ($-0.6\%$), as anchored expectations absorb much of the negative shock.}
\label{fig:missing_disinflation}
\end{figure}

Figure \ref{fig:missing_disinflation} shows the difference in outcomes: the standard NK model predicts inflation falling to $-3.7\%$, while the adaptive model reaches only $-0.6\%$. Credibility remains elevated throughout, meaning agents continue to anchor expectations at target rather than extrapolating from falling inflation. The mechanism mirrors the inflation persistence story: anchored expectations prevent spirals in both directions. This provides a structural interpretation of the anchoring explanation proposed by \citet{ball2011}. In their account, well-anchored expectations prevented deflationary spirals; our model makes this mechanism explicit and quantifiable. The much smaller deflationary response arises because high-$\theta$ agents continue forecasting 2\% inflation even as realized inflation falls, breaking the feedback loop that would otherwise amplify the demand collapse into sustained deflation. This stabilization came ``for free''; it required no additional policy action beyond the credibility inherited from previous decades.

The inflation surge of 2021--2023 provides a natural test of the two-phase dynamics. Following unprecedented fiscal and monetary stimulus combined with persistent supply chain disruptions, inflation rose from under 2\% to over 9\% by mid-2022. Initially characterized as ``transitory,'' inflation proved stickier than many forecasts suggested. Our framework interprets this through credibility dynamics: initial anchoring dampened the spike, but the delayed policy response allowed credibility to erode, generating persistence.

We simulate this episode with multiple shocks (supply chain disruptions, demand surge, and energy price increases) hitting an economy with moderate initial credibility ($\theta_0 = 0.6$) under a weak policy response ($\phi_\pi = 1.1$), capturing the Fed's delayed tightening.

\begin{figure}[htpb!]
\centering
\includegraphics[width=\textwidth]{../output/figures/section_3/figure_09_post_pandemic.png}
\caption{Post-Pandemic Inflation (2021--2024). Left panel: The adaptive model peaks at 7.8\%, while standard NK overpredicts at 16.7\%. Center panel: Credibility erodes from 0.6 to 0.10. Right panel: Two phases (light shading: dampening, dark shading: persistence), with crossover at 3.5 years.}
\label{fig:post_pandemic}
\end{figure}

Figure \ref{fig:post_pandemic} shows the two-phase dynamics. The adaptive model peaks at 7.8\%, substantially below the standard NK model's 16.7\%, illustrating the dampening role of initial anchoring. Credibility erodes from 0.6 to 0.10 as persistent inflation rewards backward-looking forecasts. This de-anchoring generates a persistence phase. Between years 3.5 and 7, the adaptive model runs approximately 1 percentage point above the NK benchmark as eroded expectations slow the return to target. The crossover at 3.5 years marks when the costs of eroded credibility begin to outweigh the initial dampening benefit.

 \paragraph{Empirical Puzzles of the New Keynesian Phillips Curve} Two well-known empirical patterns are difficult to reconcile with fixed-parameter Phillips curve specifications: monetary policy appears to operate with "long and variable" lags across different eras \citep{friedman1961lag}, and reduced-form Phillips curve estimates vary substantially across samples and time periods \citep{mavroeidis2014empirical}. Our framework suggests both patterns reflect the same underlying cause, namely time-varying credibility.

\citet{friedman1961lag} observed that monetary policy operates with ``long and variable'' lags, a pattern that challenges purely mechanical transmission models. Our framework offers a structural interpretation, since policy effectiveness varies with credibility, generating outcomes that appear as variable lags across eras.

When $\theta_t$ is high, anchored expectations do much of the stabilization automatically; agents expect inflation to return to target, and this expectation is partly self-fulfilling. When $\theta_t$ is low, backward-looking agents extrapolate recent inflation, amplifying shocks and requiring longer adjustment periods.

\begin{figure}[htpb!]
\centering
\includegraphics[width=0.7\textwidth]{../output/figures/section_3/figure_10_transmission_lags.png}
\caption{Long and variable lags. Inflation responses to identical cost-push shocks under different credibility regimes. Convergence times (marked) vary from 10 quarters under high credibility to 14 quarters under low credibility: a 40\% difference from the same shock.}
\label{fig:transmission_lags}
\end{figure}

Figure \ref{fig:transmission_lags} illustrates this mechanism. Three credibility regimes face an identical cost-push shock (0.75 percentage points quarterly, persistence $\rho = 0.75$). Higher initial credibility produces both a lower peak and faster convergence. Inflation peaks near 6\% under high credibility (circles) but close to 9\% under low credibility (squares), and convergence takes 10 quarters versus 14, a 40\% difference from the same shock.

This offers a reinterpretation of Friedman's observation: rather than viewing variable lags as evidence of inherent unpredictability in monetary transmission, our framework suggests they reflect systematic variation in the credibility state. A central bank entering a tightening cycle with high credibility can expect faster transmission than one starting from eroded credibility. The implication is that transmission speed is partly endogenous: by maintaining credibility during calm periods, central banks can ensure faster stabilization when they need it most.

\paragraph{Phillips Curve Instability} The same mechanism that generates variable transmission lags also produces apparent instability in reduced-form Phillips curve estimates. To illustrate, we simulate 200 quarters of demand shocks drawn from $N(0, 0.01^2)$ with persistence $\rho = 0.3$ under each credibility regime, holding $\theta$ fixed at its initial value. When $\theta_t$ is low, the backward-looking term dominates expectations, making inflation inherit persistence from its own past. When $\theta_t$ is high, expectations anchor to target and inflation behaves more like a jump variable. An econometrician ignoring this time-variation will find high persistence and large backward-looking coefficients in low-$\theta$ samples, and the opposite in high-$\theta$ samples.

\begin{figure}[htpb!]
\centering
\includegraphics[width=\textwidth]{../output/figures/section_3/figure_11_phillips_curve.png}
\caption{Phillips Curve instability across credibility regimes. All simulations use identical demand shocks ($\sigma = 0.01$, $\rho = 0.3$) with credibility held fixed. Panel (a): PC slope estimates from standard ($\pi_t = c + \kappa y_t$) versus hybrid ($\pi_t = c + \gamma_b \pi_{t-1} + \kappa y_t$) specifications; standard estimates vary with credibility while hybrid estimates are stable. Panel (b): AR(1) persistence.}
\label{fig:pc_instability}
\end{figure}

\paragraph{Policy Design Under Endogenous Credibility} The preceding analysis illustrates that credibility shapes inflation dynamics in systematic, predictable ways. We now turn to the normative question: how should monetary policy be designed when credibility is endogenous?

Standard New Keynesian models typically abstract from credibility dynamics. \citet{clarida1999} develop optimal policy rules under the assumption that private agents fully understand and believe the central bank's commitment; credibility is complete and immutable. The canonical framework presented in \citet{gali2015monetary}, for instance, contains no formal state variable for public beliefs about policy, implicitly treating credibility as resolved rather than evolving—a feature that limits analysis of credibility erosion during crises. More recent work has begun to address this gap: \citet{erceg2003imperfect} show how imperfect credibility generates endogenous inflation persistence and affects optimal policy design.

Our framework contributes to this literature by making credibility fully endogenous. Credibility is measured by $\theta_t$: the fraction of agents anchoring expectations to target rather than extrapolating from experience. Unlike models that assume credibility or treat it as an exogenous parameter, here credibility is earned through the central bank's track record and can erode when policy fails to deliver on its promises. This has direct implications for optimal policy, which we explore along three dimensions: (1) how aggressively should policy respond to shocks? (2) when can policymakers afford to ``wait and see,'' and when must they act immediately? (3) when credibility is low and inflationary pressure is persistent, how aggressive should disinflation be?

\paragraph{State-Dependent Optimal Policy}

The variable transmission lags documented above carry direct implications for optimal policy design. When credibility is low, stabilization takes longer and costs more; cumulative welfare losses under low credibility are nearly four times higher than under high credibility for identical shocks. This asymmetry suggests that optimal policy should itself depend on the credibility state.

To investigate this, we simulate the model across a grid of initial credibility levels ($\theta_0 \in [0.05, 0.95]$) and policy coefficients ($\phi_\pi \in [0.5, 5.0]$), subjecting each combination to the same persistent cost-push shock.\footnote{For this exercise, we implement a slight variation of the learning rule described previously. The discrete switching rule creates discontinuities in the welfare surface: small changes in $\phi_\pi$ can flip the indicator function in specific periods, causing discrete jumps in the loss. To ensure a smooth optimization problem, we replace the hard indicator with a softmax rule, where the probability of choosing each forecasting rule depends continuously on relative performance. This preserves the model's tipping-point dynamics while yielding a well-behaved objective function. See Appendix~\ref{app:computational} for details.} We evaluate performance using a standard quadratic loss function $L = \sum (\pi_t - \pi^*)^2 + 0.25 \sum y_t^2$. For each $(\theta_0, \phi_\pi)$ pair, we run a 16-quarter simulation (long enough to capture the shock response but short enough that the initial credibility state remains relevant), and compute the average squared deviations from target.

\begin{figure}[htpb!]
\centering
\includegraphics[width=\textwidth]{../output/figures/section_3/figure_12_policy_asymmetry.png}
\caption{State-dependent optimal policy. Panel (a): welfare loss as a function of $\phi_\pi$ for different initial credibility levels. All curves exhibit a U-shape: losses rise for both excessively passive ($\phi_\pi < 1$) and excessively aggressive policy, but the minimum shifts rightward as credibility falls. Panel (b): the optimal $\phi_\pi^*$ that minimizes loss at each $\theta_0$. The dashed line marks the Taylor principle threshold ($\phi_\pi = 1$).}
\label{fig:policy_asymmetry}
\end{figure}

Figure \ref{fig:policy_asymmetry}(a) displays the loss curves for four representative credibility levels. All curves exhibit a characteristic U-shape: welfare losses rise both when policy is too passive (failing to stabilize inflation) and when policy is too aggressive (inducing unnecessary output volatility). However, the location of the minimum shifts systematically with credibility. At low $\theta_0$, backward-looking expectations dominate, inflation shocks persist, and aggressive policy ($\phi_\pi \approx 3$) is needed to break the persistence. At high $\theta_0$, expectations are anchored at target, the Phillips curve effectively flattens, and moderate policy suffices; indeed, the loss curve becomes nearly flat across a wide range of $\phi_\pi$. 

Panel (b) traces out the optimal policy coefficient as a function of initial credibility. The downward slope is consistent with the intuition: central banks lacking credibility must act aggressively, while more credible ones can afford restraint. This finding aligns with \citet{nagel2024}, who shows that experience-based expectations adjust very slowly to realized inflation, implying that credibility-constrained central banks must maintain aggressive policy stances for extended periods to shift long-run beliefs toward target. Notably, the optimal $\phi_\pi^*$ falls below the Taylor principle threshold for some values of $\theta_0$. We allow this without concern for indeterminacy following \citet{angeletos2022}, who disentangle the equilibrium selection and stabilization roles of monetary policy, showing that indeterminacy is fragile to small departures from full rationality.

\paragraph{The Credibility Buffer} Central banks often advocate ``looking through'' temporary supply shocks rather than responding aggressively, a strategy that avoids unnecessary output volatility when the shock is transitory \citep[see, e.g.,][]{bernanke2007}. Our framework suggests a precondition for this patience: credibility must be sufficiently high.

To formalize this intuition, we first consider a simple test: can a high-credibility central bank simply remain passive through a supply shock? We simulate a persistent cost-push shock (5 annualized, $\rho = 0.75$) under high ($\theta_0 = 0.95$) and low ($\theta_0 = 0.15$) credibility, comparing passive policy ($\phi_\pi = 1.1$, just above the Taylor principle) against active policy ($\phi_\pi = 2.0$).

\begin{table}[htbp]
\centering
\caption{The Credibility Buffer: Riding the Wave}
\label{tab:credibility_buffer}
\begin{tabular}{lcccc}
\toprule
Scenario & Peak Inflation & Welfare Loss & Final $\theta$ & Policy Effect
\\
\midrule
High $\theta$, Passive & 7.0\% & 0.0005 & 0.895 & --- \\
High $\theta$, Active & 6.8\% & 0.0006 & 0.895 & Hurts (+28\%) \\
\\[-0.5ex]
Low $\theta$, Passive & 12.8\% & 0.0053 & 0.054 & --- \\
Low $\theta$, Active & 11.1\% & 0.0043 & 0.244 & Helps ($-19\%$) \\
\bottomrule
\end{tabular}
\begin{tablenotes}
\small
\item \textit{Notes}: Shock: 5 annualized with $\rho = 0.75$. Passive: $\phi_\pi = 1.1$. Active: $\phi_\pi = 2.0$. Policy effect shows welfare change from switching to active policy.
\end{tablenotes}
\end{table}

Table~\ref{tab:credibility_buffer} shows a clear asymmetry. Under high credibility, active policy increases welfare losses by 28\% relative to passive policy. The inflation reduction is marginal (7.0 to 6.8 peak), while the output costs are nontrivial. Anchored expectations absorb the shock; aggressive intervention adds little on the inflation margin while incurring output costs. Credibility barely erodes, falling only from 0.95 to 0.90. Under low credibility, the pattern reverses: passive policy allows inflation to reach 12.8 and credibility to collapse from 0.15 to 0.054. Active policy reduces welfare losses by 19\%.

The preceding analysis assumes the central bank knows the shock is transitory. In practice, central bankers face uncertainty, as a shock that appears transitory may prove persistent. This raises a subtler question. Can a credible central bank afford to wait and see before responding, buying time to assess the shock's nature?

To address this, we simulate identical persistent supply shocks under high ($\theta_0 = 0.95$) and low ($\theta_0 = 0.50$) credibility, comparing two policy responses: immediate, where the Taylor rule is active from $t=0$, and delayed, where the nominal interest rate is held fixed at its neutral level for 8 quarters before the Taylor rule engages.

\begin{figure}[htpb!]
\centering
\includegraphics[width=\textwidth]{../output/figures/section_3/figure_13_credibility_buffer.png}
\caption{The option to wait depends on credibility. Panel (a): inflation paths under immediate (solid) versus delayed (dashed) policy response for high credibility (blue) and low credibility (black). The shaded region marks the delay period. Panel (b): cumulative welfare loss for each scenario.}
\label{fig:credibility_buffer}
\end{figure}

Figure~\ref{fig:credibility_buffer} shows the results. Under high credibility, the delayed response tracks the immediate response closely, as anchored expectations absorb the shock even during the delay period. Under low credibility, however, the delayed response generates higher and more persistent inflation. The mechanism is the de-anchoring spiral: each period of elevated inflation feeds into higher expected inflation, which generates higher actual inflation. Panel (b) quantifies the cost: welfare losses under delayed response are similar to immediate response when credibility is high, but diverge substantially when credibility is low.

Together, these experiments suggest that the credibility buffer operates at two levels. First, for shocks that are clearly transitory, high-credibility central banks can remain passive throughout, as active policy adds output costs with little inflation benefit. Second, even when the shock's persistence is uncertain, high-credibility central banks can afford to wait before committing to a response. A central bank lacking credibility cannot afford this patience: it must respond both actively and promptly to prevent expectation de-anchoring.

\paragraph{Regime Change: Shock Therapy vs Gradualism}

The credibility buffer analysis suggests that low-credibility central banks cannot afford to be passive. But this raises a follow-up question. Given that action is required, how aggressive should it be? Should a central bank facing a credibility crisis act aggressively from day one (``shock therapy''), or should policy tighten gradually? This debate has historical precedent in the contrast between the decisive 1979 tightening and the hesitant policy of the 1970s \citep[see][]{goodfriend2005}.

Our framework suggests shock therapy tends to outperform gradualism when initial conditions are adverse. The mechanism operates through credibility dynamics. When policy is aggressive from day one, inflation falls quickly, pulling down backward-looking forecasts. The CB-anchored rule begins outperforming sooner, and credibility recovers earlier. Under gradualism, aggressive policy arrives only after a ramp-up period. During this transition, inflation remains elevated, backward-looking forecasts stay high, and the CB-anchored rule continues losing the performance comparison. By the time policy reaches full strength, credibility may have eroded further.

\begin{table}[htbp]
\centering
\caption{Shock Therapy vs Gradualism in a Credibility Crisis}
\label{tab:shock_gradual}
\begin{tabular}{lccc}
\toprule
Strategy & Welfare Loss & Recovery Time & Final $\theta$ \\
& & (quarters to $\theta > 0.5$) & \\
\midrule
Shock therapy & 0.0475 & 44 & 0.914 \\
Gradualism & 0.0480 & 48 & 0.908 \\
Constant moderate & 0.0624 & Never & 0.010 \\
\bottomrule
\end{tabular}
\begin{tablenotes}
\small
\item \textit{Notes}: Initial conditions: $\pi_0 = 14$, $\theta_0 = 0.05$. Shock therapy: $\phi_\pi = 3.5$ throughout. Gradualism: $\phi_\pi$ ramps from 1.5 to 3.5 over 20 quarters. Constant moderate: $\phi_\pi = 1.5$ throughout. Ongoing shocks: 2.4 with $\rho = 0.8$.
\end{tablenotes}
\end{table}

Table~\ref{tab:shock_gradual} presents simulation results for a crisis scenario: 14 initial inflation and low credibility ($\theta_0 = 0.05$). Shock therapy outperforms gradualism on both welfare and recovery speed, though the differences are modest (1 welfare improvement, 4 quarters faster recovery). The more notable result is the constant moderate baseline: with $\phi_\pi = 1.5$, credibility falls to 0.01 and does not recover within the 20-year horizon. When credibility is low and inflationary pressure is persistent, insufficiently aggressive policy risks a self-reinforcing de-anchoring dynamic.

These results connect to historical experience. The intermittent policies of the 1970s can be characterized as responding to inflation, but without the sustained commitment needed to shift backward-looking dynamics \citep{goodfriend2005}. Each temporary tightening was followed by easing before credibility could recover. The 1979 policy shift took a different approach: aggressive and sustained. The 1981–82 recession was severe, but credibility recovered within a few years.

%%%%%%%%%%%%%%%%%%%%%%%%%%%%%%%%%%%%%%%%%%%%%%%%%%%%%%%%%%%%%%%%%%%%%%%%%%%%%%%%%%%%%%%%%%%%%%%%%%%%%%%%%%%%%%%%%%%%%%%%%%%%%%%%%%%%%%%%%%%%%%%%%%%%%%%%%%%%%%%%%%%%%%%%%%%%%%%%%%%%%%%%%%%%%%%%%%%%%%%%%%%%%%%%%%%%%%%%%%%%%%%%%%%%%%%%%%%%%%%%%%%%%%%%%%%%%%%%%%%%%%%%%%%%%%%%%%%%%%%%%%%%%%%%%%%%%%%%%%%%%%%%%%%%%%%%%%%%%%%%%%%%%%%%%%%%%%%%%%%%%%%%%%%%%%%%%%%%%%%%%%%%%%%%%%%%%%%%%%%%%%%%%%%%%%%%%%%%%%%%%%%%%%%%%%%%%%%%%%%%%%%%%%%%%%%%%%%%%%%%%%%%%%%%%%%%%%%%%%%%%%%%%%%%%%%%%%%%%%%%%%%%%%%%%%%%%%%%%%%%%%%%%%%%%%%%%%%%%%%%%%%%%%%%%%%%%%%%%%%%%%%%%%%%%%%%%%%%%%%%%%%%%%%%%%%%%%%%%%%%%%%%%%%%%%%%%%%%%%%%%%%%%%%%%%%%%%%%%%%%%%%%%%%%%%%%%%%%%%%%%%%%%%%%%%%%%%%%%%%%%%%%%%%%%%%%%%%%%%%%%%%%%%%%%%%%%%%%%%%%%%%%%%%%%%%%%%%%%%%%%%%%%%%%%%%%%%%%%%%%%%%%%%%%%%

\section{The MAB Framework: Extensions and Variations}
\label{sec:extensions}

The specific choices in the preceding sections - two forecasting rules, a fixed learning rate, discrete switching - represent one instance of a more general approach. The underlying logic is portable, as agents draw on multiple information sources, evaluate them based on forecasting performance, and shift toward whichever rule has worked, all without requiring knowledge of the economy's structure.

We present two extensions that illustrate this flexibility. First, we introduce trend-following expectations, in which agents extrapolate not just the level but the direction of recent inflation. This third forecasting rule allows the framework to speak to hyperinflationary spirals, where expectations and inflation feed on each other explosively, and to the dynamics of stabilization programs like Brazil's Real Plan. Second, we replace the fixed memory window with exponentially discounted evaluation of past performance, capturing how inflation trauma can leave lasting scars on expectations. This extension addresses persistent deflation traps like Japan's, and why countries that experienced comparable inflationary crises can exhibit very different credibility dynamics depending on how recently those crises occurred. We do not claim these extensions ``solve'' these puzzles, but they offer a lens through which to interpret them, one grounded in the same adaptive logic developed above.

\paragraph{The Three-Arm Model} Our first extension illustrates the gear-shift logic previewed in the introduction: when existing heuristics fail, agents shift to more complex forecasting rules. When inflation is stable, trusting the central bank's target works well. When inflation deviates, agents may shift to tracking recent inflation. But when inflation is accelerating, even backward-looking forecasts systematically underpredict. At that point, agents may extrapolate the trend: if inflation rose from 5\% to 7\%, forecast 9\%. This logic echoes what Brazilian economists termed "inertial inflation" \citep{bresserpereira2023inertial}: when agents index expectations to recent price movements, inflation acquires a momentum of its own that persists even as the original shocks fade.

The trend-following (TF) rule takes the form $\mathbb{E}[\pi] = \pi_{t-1} + (\pi_{t-1} - \pi_{t-2})$: last period's inflation plus the recent change. Survey evidence supports such behavior. \cite{coibion2015} document that household inflation expectations exhibit substantial extrapolation. \cite{malmendier2016learning} show that personal inflation experiences shape expectations for decades, with agents overweighting recent trends. In experimental settings, \cite{hommes2017} find that subjects spontaneously adopt trend-following rules when prices exhibit momentum.

The selection mechanism extends naturally to three rules. Agents evaluate forecasting performance as before, with aggregate expectations given by:

\begin{equation}
\mathbb{E}_t[\pi_{t+1}] = \theta^{CB}_t \cdot \pi^* + \theta^{BL}_t \cdot \pi_{t-1} + \theta^{TF}_t \cdot \left[\pi_{t-1} + (\pi_{t-1} -\pi_{t-2})\right]
\end{equation}

where $\theta^{CB}_t + \theta^{BL}_t + \theta^{TF}_t = 1$. As in the baseline, agents prefer simpler rules when forecasting performance is similar: a complexity hierarchy $\varepsilon_{CB} < \varepsilon_{BL} <\varepsilon_{TF}$ ensures trend-following remains dormant during normal times.

\paragraph{Hyperinflation and the Gear Shift} The two-arm model generates persistence but not acceleration—the explosive dynamics characteristic of hyperinflation. The three-arm model fills this gap.

Trend-following outperforms backward-looking only when inflation is accelerating. If inflation is high but stable, backward-looking predicts accurately while trend-following over-predicts. The ``gear shift'' to trend-following thus requires sustained acceleration—which itself requires monetary policy too passive to stabilize. Under active policy, inflation may rise but eventually stabilizes; backward-looking captures the dynamics and trend-following never activates. Under sufficiently passive policy, inflation can accelerate continuously, trend-following begins outperforming, and extrapolation creates self-fulfilling acceleration.

The path to hyperinflation is thus a two-stage de-anchoring: first CB loses to BL (agents stop trusting the target), then BL loses to TF (agents start extrapolating trends). This suggests that preventing the shift to trend-following requires early action. Once TF dominates expectations, stabilization becomes costly because any slowdown in inflation must persist long enough for simpler rules to regain their performance advantage.

\paragraph{Brazil and the Real Plan} The three-arm model offers a perspective on Brazil's Real Plan (1994). After years of failed stabilization attempts, inflation fell from approximately 50 per month to single digits when the Real was introduced in July 1994. Our framework suggests the key lies not in the policy rule itself, but in the institutional innovation that preceded it: the Unidade Real de Valor (URV).

From the late 1980s through early 1994, Brazil experienced sustained hyperinflation with monthly rates often exceeding 40. Previous stabilization attempts - the Cruzado Plan (1986), Bresser Plan (1987), Summer Plan (1989), and Collor Plans (1990–91) - all failed. Each introduced price freezes or new currencies directly into an economy where trend-following expectations dominated. Our framework suggests why: when agents expect inflation to accelerate, even small shocks get amplified through extrapolative forecasts, overwhelming whatever stabilization policy is in place.

The Real Plan took a different approach. Rather than attempting to stabilize the existing currency directly, it first introduced the URV, a unit of account pegged to the US dollar. Prices and wages were quoted in URV while the old Cruzeiro Real continued to circulate for transactions. Because of the peg, URV-denominated prices were stable by construction. After several months, once URV pricing had become widespread, the Real replaced the URV at parity. Monthly inflation fell from over 45 in June 1994 to under 2 by August.

The key to modeling this episode is that the URV does not start with fresh expectations. It inherits the de-anchored state from the Cruzeiro era, i.e, agents initially remain dominated by trend-following ($\theta_{TF} = 0.80$, $\theta_{CB} = 0.05$). What changes is the inflation they observe: URV-denominated prices are stable. As the MAB learning runs, trend-following systematically over-predicts (it expects continued acceleration), while the CB-anchored rule is accurate. Expectations re-anchor not because agents are told to trust the new system, but because stable prices make the CB-anchored forecast perform well.

We formalize this by comparing two scenarios facing identical cost-push shocks (2\% for four quarters) under identical active monetary policy ($\phi_\pi = 1.5$). The counterfactual starts from the inherited hyperinflation expectation state ($\theta_{TF} = 0.80$, $\theta_{CB} = 0.05$): an economy where 80\% of agents are already extrapolating trends and only 5\% trust the central bank. The Real Plan scenario first simulates a four-quarter URV phase, during which expectations evolve via MAB learning while inflation is held at target, then runs the NK model from the resulting---substantially re-anchored---expectation state.

\begin{figure}[htpb!]
\centering
\includegraphics[width=\textwidth]{../output/figures/section_4/figure_15_brazil.png}
\caption{Brazil: Counterfactual vs Real Plan. Both simulations face identical shocks under $\phi_\pi = 1.5$. Panel~(a): Inflation paths. Panel~(b): Expectation shares during the URV phase. Panel~(c): Credibility dynamics during stabilization.}
\label{fig:brazil}
\end{figure}

Figure~\ref{fig:brazil} shows the results. Panel~(a) displays inflation paths: the counterfactual produces peak inflation exceeding 2000\%, while the Real Plan path peaks at 18\% before returning to target. The 2\% shock is a trigger, not the driver: the explosive dynamics arise because the counterfactual economy is a powder keg, with 80\% of agents feeding trend-following forecasts into an accelerating spiral. The Real Plan comparison works precisely because the URV phase drains this expectation fuel before the spark arrives. Both scenarios face identical shocks under identical policy; the only difference is the composition of expectations inherited at the start.

Panel~(b) illustrates the URV mechanism. Starting from $\theta_{TF} = 0.80$, we hold inflation at target and track how expectations evolve. Within four quarters, CB-anchored expectations rise from 5 to 41. Trend-following collapses because it over-predicts when inflation is stable; the CB-anchored rule, being exactly correct, accumulates performance advantage.

Panel~(c) shows credibility dynamics during stabilization. In the counterfactual, credibility falls toward zero. Under Real Plan conditions, credibility rises steadily as smooth disinflation rewards the CB-anchored forecast. Even accounting for the cost of continued Cruzeiro hyperinflation during the URV period, the Real Plan path reduces welfare losses by over 90\%.

Argentina's Convertibility Plan (1991) followed similar logic: a currency board pegged one-to-one to the US dollar provided a credible nominal anchor, and inflation fell from over 2000 in 1990 to single digits by 1993. The broader lesson is that direct stabilization from hyperinflation may fail not because of insufficient policy commitment, but because de-anchored expectations amplify any remaining shocks. A nominal anchor can break this cycle by providing a period of observed stability during which expectations re-anchor through learning.

\paragraph{Long Memory, Trauma, and Skittish Expectations}

The baseline model evaluates forecasting performance over a fixed $k$-period window, weighting all observations equally. This is a convenient simplification, but it abstracts from an important feature of credibility dynamics emphasized in the monetary policy literature: credibility is earned gradually over long horizons. \citet{goodfriend2005} document that markets took years to believe the Volcker disinflation would persist, as long-term interest rates remained elevated well after inflation had fallen, reflecting the slow accumulation of trust in the Fed's commitment. More broadly, \citet{bordosiklos2015} characterize credibility as something that ``evolves in a non-linear manner, i.e.\ it is earned slowly and painstakingly yet is susceptible to evaporate on a moment's notice.'' A three-quarter memory window cannot capture this gradual building of credibility capital.

The lasting scars that inflation experiences leave on expectations present a similar challenge. \citet{malmendier2016learning} document that personal inflation experiences shape expectations with a half-life of several years: Germans who lived through the Weimar hyperinflation remained inflation-averse for decades, and Americans who came of age during the 1970s respond differently to price signals than those raised during the Great Moderation. A fixed memory window, which treats all recent observations equally, cannot capture this generational imprinting.

To capture these dynamics, we replace the fixed window with exponentially discounted losses. Rather than evaluating each forecasting rule over the past $k$ periods with equal weight, agents now discount distant observations:

\begin{equation}
L^j_t = \sum_{s=0}^{\infty} \delta^s \left(\pi_{t-s} -
\mathbb{E}^j_{t-s-1}[\pi_{t-s}]\right)^2
\end{equation}

where $\delta \in (0,1)$ is a discount factor. Higher $\delta$ means longer memory: agents weight distant experiences more heavily relative to recent ones. With $\delta = 0.8$, the half-life is approximately 3 quarters; recent observations dominate, but the past is not forgotten.

This formulation captures how people who lived through high inflation remain more responsive to inflation signals and quicker to abandon the central bank anchor. In our framework, this manifests as country-specific dynamics depending on inflation history.
  
We illustrate by comparing three hypothetical countries with different 20-year histories, varying both the severity and timing of inflation episodes. The first is a country with stable inflation at target throughout. The second experienced hyperinflation (50\% monthly) that ended 15 years ago---enough time for partial but not complete recovery. The third has ongoing chronic moderate inflation (7.5\% annually) for the past decade. The question is how these different histories shape the response to a new shock today.

We first run each country through its 20-year history to establish its credibility state, then hit all three with an identical cost-push shock. Figure~\ref{fig:long_memory} shows the results.

\begin{figure}[htpb!]
\centering
\includegraphics[width=\textwidth]{../output/figures/section_4/figure_16_long_memory.png}
\caption{Long Memory: Same Shock, Different Outcomes ($\delta = 0.8$). All countries face an identical cost-push shock. Panel (a): Inflation response. Panel (b): Credibility dynamics. The stable country maintains high credibility; distant trauma has partially recovered ($\theta \approx 0.5$); chronic current inflation leaves credibility near zero.}
\label{fig:long_memory}
\end{figure}

Panel (a) shows inflation responses. The stable country, with full credibility ($\theta = 1$), sees peak inflation of 8\% before returning smoothly to target. The distant trauma country, despite 15 years of stability since its hyperinflation, has only partially recovered ($\theta \approx 0.5$) and experiences a larger response: peak inflation of 12\%. The chronic current country, with credibility near zero ($\theta \approx 0$), sees the shock amplified substantially to 23\%. The same shock produces outcomes that differ by a factor of three, purely because of inherited credibility.

Panel (b) shows credibility dynamics during the shock. The stable country maintains high credibility throughout. The distant trauma country, starting from intermediate credibility, actually sees credibility \textit{rise} as the shock passes and inflation returns to target. The chronic current country sees credibility collapse further---the shock reinforces its de-anchored state rather than being absorbed.

The distant trauma case is particularly instructive, as even 15 years after a hyperinflation ended, credibility has only partially recovered. This slow rebuilding of trust, combined with the vulnerability to new shocks, offers a perspective on why countries with inflation histories remain more susceptible to de-anchoring episodes decades later.

\paragraph{Japan: Expectations Stuck Below Target} The learning mechanism cuts both ways: a central bank that consistently undershoots its target loses credibility just as surely as one that overshoots. Japan since the 1990s illustrates this. For over two decades, Japan experienced near-zero inflation despite an official 2\% target and unprecedented monetary easing, with inflation expectations stubbornly stuck below target, a puzzle that likely involves multiple factors, including the zero lower bound. Our framework suggests one contributing mechanism: when the central bank promises 2\% but delivers 0\% year after year, the CB-anchored forecasting rule accumulates systematic errors. Agents learn that trusting the target is a losing strategy, and expectations become anchored, but at the wrong level.

Simulations confirm this intuition: if inflation stays persistently at 0\% while the target is 2\%, credibility collapses just as it would with persistent 4\% inflation (see Online Appendix for details). Small misses, say, 1\% versus 2\%, fall within the tie-breaking threshold and do not trigger de-anchoring, but a 2 percentage point persistent miss is large enough that the CB-anchored rule loses. The policy implication is that escaping a below-target trap may be as difficult as escaping an above-target one, because in both cases the central bank must overcome a stock of accumulated credibility losses, and the same anchoring mechanism that stabilizes inflation near target when credibility is well-placed works against the central bank when expectations have anchored at the wrong level.


%%%%%%%%%%%%%%%%%%%%%%%%%%%%%%%%%%%%%%%%%%%%%%%%%%%%%%%%%%%%%%%%%%%%%%%%%%%%%%%%%%%%%%%%%%%%%%%%%%%%%%%%%%%%%%%%%%%%%%%%%%%%%%%%%%%%%%%%%%%%%%%%%%%%%%%%%%%%%%%%%%%%%%%%%%%%%%%%%%%%%%%%%%%%%%%%%%%%%%%%%%%%%%%%%%%%%%%%%%%%%%%%%%%%%%%%%%%%%%%%%%%%%%%%%%%%%%%%%%%%%%%%%%%%%%%%%%%%%%%%%%%%%%%%%%%%%%%%%%%%%%%%%%%%%%%%%%%%%%%%%%%%%%%%%%%%%%%%%%%%%%%%%%%%%%%%%%%%%%%%%%%%%%%%%%%%%%%%%%%%%%%%%%%%%%%%%%%%%%%%%%%%%%%%%%%%%%%%%%%%%%%%%%%%%%%%%%%%%%%%%%%%%%%%%%%%%%%%%%%%%%%%%%%%%%%%%%%%%%%%%%%%%%%%%%%%%%%%%%%%%%%%%%%%%%%%%%%%%%%%%%%%%%%%%%%%%%%%%%%%%%%%%%%%%%%%%%%%%%%%%%%%%%%%%%%%%%%%%%%%%%%%%%%%%%%%%%%%%%%%%%%%%%%%%%%%%%%%%%%%%%%%%%%%%%%%%%%%%%%%%%%%%%%%%%%%%%%%%%%%%%%%%%%%%%%%%%%%%%%%%%%%%%%%%%%%%%%%%%%%%%%%%%%%%%%%%%%%%%%%%%%%%%%%%%%%%%%%%%%%%%%%%%%%%%%%%%%



\section{Conclusion}

This paper develops an adaptive model of inflation expectations in which agents dynamically choose between anchoring to the central bank's target and extrapolating from recent inflation, based on forecasting performance. Central bank credibility, measured by the fraction of agents trusting the target, emerges endogenously from this process: it is earned through policy performance and eroded by persistent deviations.

\paragraph{Main Findings} The framework provides a lens for understanding several phenomena that standard models struggle to reconcile. The 1970s oil shocks and the 2008 commodity price surge were similar in magnitude but produced substantially different outcomes; in our framework, this reflects differences in inherited credibility rather than structural change. The Great Moderation emerges as a period where high credibility acted as an automatic stabilizer; the ``missing disinflation'' of 2008--2012 reflects anchored expectations absorbing the demand collapse; and the post-COVID inflation surge illustrates two-phase dynamics, with initial dampening from residual anchoring followed by persistence as credibility eroded. The time-varying inflation persistence documented in U.S. data aligns with the model's central prediction: persistence rises when credibility falls, and falls when credibility rises.

\paragraph{Extensions} The framework accommodates alternative specifications that address more extreme phenomena. Adding trend-following expectations (a third forecasting rule activated when inflation accelerates) allows the model to speak to hyperinflationary spirals and stabilization programs like Brazil's Real Plan, where the URV re-anchored expectations before the new currency launched. Replacing the fixed memory window with exponentially discounted evaluation captures how inflation trauma leaves lasting scars, offering a perspective on Japan's deflation trap and cross-country differences in credibility dynamics.

\paragraph{Policy Implications} The framework suggests that optimal policy is state-dependent, as credibility-constrained central banks require more aggressive responses than credible ones facing identical shocks. High credibility provides a buffer, allowing policymakers to ``look through'' temporary disturbances; low credibility eliminates this option. The model also highlights a fundamental asymmetry: credibility erodes quickly but rebuilds slowly. The costs of allowing expectations to de-anchor extend far beyond the initial shock period.

\paragraph{Limitations and Future Work} Several directions remain for future research. We assume a constant learning rate, but attention is plausibly state-dependent, as agents may learn faster during volatile periods. The paradox of faster learning (Table~\ref{tab:paradox_learning}) suggests this matters because a critical $\eta$ threshold separates recoverable de-anchoring from self-sustaining spirals, and state-dependent learning rates would shift this threshold endogenously with the inflation environment. The self-reinforcing nature of credibility dynamics suggests the model may admit multiple steady states; characterizing their stability would complement our simulation-based results. Replacing greedy rule selection with softmax (probabilistic) switching would generate persistent disagreement among agents even when they observe the same
   data, providing a more natural account of the cross-sectional heterogeneity visible in inflation expectation surveys. Finally, directly estimating the credibility state $\theta_t$ from survey data and testing whether it explains cross-sectional expectation dispersion would be an important validation.

\newpage

\bibliographystyle{ecta}
\bibliography{references}

\newpage
\appendix

%%%%%%%%%%%%%%%%%%%%%%%%%%%%%%%%%%%%%%%%%%%%%%%%%%%%%%%%%%%%%%%%%%%%%%%%%%%%%%%%%%%%%%%%%%%%%%%%%%%%%%%%%%%%%%%%%%%%%%%%%%%%%%%%%%%%%%%%%%%%%%%%%%%%%%%%%%%%%%%%%%%%%%%%%%%%%%%%%%%%%%%%%%%%%%%%%%%%%%%%%%%%%%%%%%%%%%%%%%%%%%%%%%%%%%%%%%%%%%%%%%%%%%%%%%%%%%%%%%%%%%%%%%%%%%%%%%%%%%%%%%%%%%%%%%%%%%%%%%%%%%%%%%%%%%%%%%%%%%%%%%%%%%%%%%%%%%%%%%%%%%%%%%%%%%%%%%%%%%%%%%%%%%%%%%%%%%%%%%%%%%%%%%%%%%%%%%%%%%%%%%%%%%%%%%%%%%%%%%%%%%%%%%%%%%%%%%%%%%%%%%%%%%%%%%%%%%%%%%%%%%%%%%%%%%%%%%%%%%%%%%%%%%%%%%%%%%%%%%%%%%%%%%%%%%%%%%%%%%%%%%%%%%%%%%%%%%%%%%%%%%%%%%%%%%%%%%%%%%%%%%%%%%%%%%%%%%%%%%%%%%%%%%%%%%%%%%%%%%%%%%%%%%%%%%%%%%%%%%%%%%%%%%%%%%%%%%%%%%%%%%%%%%%%%%%%%%%%%%%%%%%%%%%%%%%%%%%%%%%%%%%%%%%%%%%%%%%%%%%%%%%%%%%%%%%%%%%%%%%%%%%%%%%%%%%%%%%%%%%%%%%%%%%%%%%%%%%

\section{The Multi-Armed Bandit Algorithm}
  \label{app:mab}

This appendix provides background on the Multi-Armed Bandit (MAB) framework that inspired our expectation formation process. The goal is accessibility rather than comprehensiveness: we introduce the key concepts economists need to understand our modeling choices, and point readers toward more complete treatments. For a rigorous introduction to MAB and reinforcement learning more broadly, see \citet{sutton2018reinforcement}; for evidence that these algorithms describe human decision-making in cognitive tasks, see \citet{binz2023}.

\paragraph{The Problem} The MAB problem takes its name from a gambler facing several slot machines (colloquially called ``one-armed bandits''), each with an unknown probability of paying out. The gambler must decide which machine to play on each turn to maximize cumulative winnings. The challenge is the explore-exploit tradeoff: should the gambler exploit the machine that has paid out most so far, or explore other machines that might be better but are less well understood?

This tradeoff captures the essence of forecasting rule choice in our model. Agents must decide whether to rely on a rule that has worked well recently (exploit) or consider alternatives that might perform better under current conditions (explore). Unlike settings where the optimal choice is stable, macroeconomic regimes shift: a rule that worked during the Great Moderation may fail during an inflation surge, and vice versa. MAB algorithms are designed precisely for such non-stationary environments where past performance provides useful but imperfect guidance about future success. Crucially, this learning process requires no understanding of monetary policy, Phillips curves, or general equilibrium: agents need only observe which forecasting rule has predicted inflation more accurately, a comparison they can make using publicly available data.

\paragraph{Algorithms and Our Implementation} Several algorithms have been developed to solve the MAB problem, differing in how they balance exploration and exploitation. The simplest is \textit{greedy} selection: always choose whichever option has the highest cumulative performance. This approach is optimal if the agent already knows which option is best, but it never explores: if early experience happens to favor a suboptimal choice, the agent may stick with it indefinitely. A smoother alternative is \textit{softmax} (or Boltzmann) selection, where the probability of choosing each option depends continuously on its estimated performance rather than switching discretely. Other approaches exist in the literature (notably $\varepsilon$-greedy and upper confidence bound methods), but we do not use them here.

Our model uses greedy selection with Calvo-style updating: each period, a fraction $\eta$ of agents reconsiders their forecasting rule and switches to whichever performed better; the remaining $(1-\eta)$ maintain their current choice. This yields:
  
\begin{equation}
\theta_{t+1} = (1-\eta)\theta_t + \eta \cdot \mathds{1}[L^{CB}_t \leq L^{BL}_t + \varepsilon]
\end{equation}

where $\theta_t$ is the fraction using the CB-anchored rule, $\mathds{1}[\cdot]$ indicates whether the anchored rule outperformed, and $\varepsilon > 0$ is a tie-breaking threshold favoring the simpler rule when losses are similar.

We choose this deliberately simple specification (greedy selection with no explicit exploration) for three reasons. First, \textit{transparency}: the mechanism is easy to understand and analyze. Second, \textit{tractability}: credibility reduces to a single state variable with a deterministic law of motion, making the model easy to embed in standard New Keynesian frameworks. Third, \textit{endogenous exploration}: even without explicit exploration, macroeconomic fluctuations expose both rules to varied conditions over time. When inflation is volatile, both rules accumulate errors, preventing either from dominating permanently, as the economy itself provides the variation that formal MAB algorithms build in by design.

The Calvo-style updating also has economic motivation. Not all agents can or do reconsider their forecasting approach each period. Inattention, information costs, and decision inertia (all well-documented in the behavioral economics literature) imply that belief revision is lumpy rather than continuous. The parameter $\eta$ captures the average frequency of reconsideration, analogous to how the Calvo parameter captures the frequency of price adjustment.

\paragraph{When We Use Softmax} For most applications, we use the hard indicator described above. However, the optimal policy analysis in Section~\ref{sec:applications} requires a smooth objective function, so we employ a softmax variant that closely approximates the hard indicator while remaining differentiable. Appendix~\ref{app:computational} provides details.

\paragraph{Relation to Adaptive Learning} Our MAB approach differs from standard adaptive learning models in a fundamental way. The literature following \citet{evans2001learning} models agents who continuously update the \textit{parameters} of a single forecasting model through recursive least squares or similar algorithms. Beliefs evolve smoothly as coefficients gradually converge toward rational expectations values.

Our framework instead models discrete switching between qualitatively different forecasting \textit{rules}. Agents choose which model to use, not what parameters to use within a model. This generates sharper regime transitions: when inflation rises persistently, many agents can simultaneously abandon the anchor, producing the rapid de-anchoring observed in historical episodes. It also generates asymmetry because switching away from a failing strategy is fast (one bad period can trigger reconsideration), while rebuilding trust in the anchor requires sustained evidence of good performance, since forecast errors are non-negative, meaning that accurate predictions do not offset past misses, which persist in the evaluation window until they age out.

Experimental evidence supports this discrete-choice structure. In repeated decision tasks, people tend to maintain a small set of strategies and switch between them based on recent performance, rather than continuously updating a single model \citep[see][for a synthesis]{binz2023}. This is precisely the behavior our framework assumes: agents do not learn the structure of the economy, they simply track which forecasting rule has been more accurate and act accordingly.

\paragraph{Connection to Bayesian Learning} MAB learning can be understood as a tractable approximation to Bayesian inference. A Bayesian agent facing the same problem would maintain a posterior distribution over the relative accuracy of each forecasting rule, updating it each period as new forecast errors are observed. Our greedy MAB implementation tracks only cumulative performance rather than a full posterior, reducing cognitive demands while preserving the essential dynamics: well-performing rules gain adherents, poorly-performing ones are abandoned.

The Bayesian formulation does offer one advantage: because agents hold probabilistic beliefs, their rule choices are stochastic, generating cross-sectional heterogeneity in expectations even among agents with identical information. This is a natural source of forecast disagreement that the deterministic MAB version lacks. Appendix~\ref{app:bayesian} formalizes this connection, presenting a Bayesian signal extraction framework as an alternative microfoundation and showing that it generates aggregate dynamics closely resembling those of the MAB model.

\paragraph{Further Reading} For readers interested in exploring these ideas further: \citet{sutton2018reinforcement} provide a comprehensive introduction to reinforcement learning, of which MAB is the simplest case (Chapters 2--3 cover bandit algorithms in detail); \citet{binz2023} show how reinforcement learning algorithms describe human cognition across a wide range of tasks, providing behavioral foundations for our approach; and \citet{charpentier2021} survey applications of reinforcement learning in economics, including connections to adaptive expectations and learning in games.

%%%%%%%%%%%%%%%%%%%%%%%%%%%%%%%%%%%%%%%%%%%%%%%%%%%%%%%%%%%%%%%%%%%%%%%%%%%%%%%%%%%%%%%%%%%%%%%%%%%%%%%%%%%%%%%%%%%%%%%%%%%%%%%%%%%%%%%%%%%%%%%%%%%%%%%%%%%%%%%%%%%%%%%%%%%%%%%%%%%%%%%%%%%%%%%%%%%%%%%%%%%%%%%%%%%%%%%%%%%%%%%%%%%%%%%%%%%%%%%%%%%%%%%%%%%%%%%%%%%%%%%%%%%%%%%%%%%%%%%%%%%%%%%%%%%%%%%%%%%%%%%%%%%%%%%%%%%%%%%%%%%%%%%%%%%%%%%%%%%%%%%%%%%%%%%%%%%%%%%%%%%%%%%%%%%%%%%%%%%%%%%%%%%%%%%%%%%%%%%%%%%%%%%%%%%%%%%%%%%%%%%%%%%%%%%%%%%%%%%%%%%%%%%%%%%%%%%%%%%%%%%%%%%%%%%%%%%%%%%%%%%%%%%%%%%%%%%%%%%%%%%%%%%%%%%%%%%%%%%%%%%%%%%%%%%%%%%%%%%%%%%%%%%%%%%%%%%%%%%%%%%%%%%%%%%%%%%%%%%%%%%%%%%%%%%%%%%%%%%%%%%%%%%%%%%%%%%%%%%%%%%%%%%%%%%%%%%%%%%%%%%%%%%%%%%%%%%%%%%%%%%%%%%%%%%%%%%%%%%%%%%%%%%%%%%%%%%%%%%%%%%%%%%%%%%%%%%%%%%%%%%%%%%%%%%%%%%%%%%%%%%%%%%%%%%%%%%

\section{A Rational Inattention Foundation for the Simplicity Bias}
\label{app:entropy}

This appendix provides a decision-theoretic foundation for the tie-breaking parameter $\varepsilon$ in the learning rule. Recall that agents switch to the CB-anchored forecast when $L^{CB}_t \leq L^{BL}_t + \varepsilon$. We first argue why this parameter is essential for the model's dynamics, then provide a back-of-the-envelope calibration grounded in the rational inattention literature.

\paragraph{The Re-Anchoring Problem} Consider an agent who begins fully anchored, forecasting $\pi^e_t = \pi^*$. A cost-push shock pushes inflation above target, and the backward-looking rule (which would have forecast last period's elevated inflation) outperforms the CB-anchored rule. The agent switches to backward-looking expectations. Now suppose monetary policy successfully brings inflation back to target gradually. At this point, with $\pi_t = \pi^*$, both rules generate identical forecasts: the CB rule forecasts $\pi^*$ by construction, while the backward-looking rule forecasts $\pi_{t-1} = \pi^*$, so both generate identical losses. Without $\varepsilon > 0$, the agent has no reason to switch back, since the CB rule offers no performance advantage. The agent remains de-anchored despite the central bank's success. This is the re-anchoring problem: once credibility is lost, it cannot be regained through good policy alone.

This is not a minor issue but a major failure to match historical experience. Central banks have repeatedly restored credibility following periods of elevated inflation: the Volcker disinflation being the canonical example. A model that traps the economy in permanent de-anchoring after any sufficiently large shock misses a central feature of monetary history.

Figure~\ref{fig:reanchoring_problem} illustrates this problem. We subject the economy to a large persistent cost-push shock and compare two parameterizations: $\varepsilon = 0$ and $\varepsilon = 10^{-4}$. In both cases, the shock triggers identical de-anchoring dynamics, and monetary policy eventually stabilizes inflation at the 2\% target. With $\varepsilon = 0$, credibility remains near zero indefinitely: agents have no incentive to switch back because both rules now yield identical predictions. With $\varepsilon = 10^{-4}$, credibility recovers once inflation stabilizes: the CB-anchored rule, being simpler, is preferred when performance is equivalent.

\begin{figure}[htbp]
\centering
\includegraphics[width=\textwidth, height = \textheight, keepaspectratio]{../output/figures/appendix_B/figure_B1_reanchoring.png}
\caption{The re-anchoring problem. Both columns show the response to a large persistent cost-push shock (2\%, $\rho = 0.8$) with baseline learning speed $\eta = 0.10$. Left column: without tie-breaking ($\varepsilon = 0$), credibility cannot recover even after inflation returns to target. Right column: with tie-breaking ($\varepsilon = 10^{-4}$), credibility recovers once inflation stabilizes.}
\label{fig:reanchoring_problem}
\end{figure}

\paragraph{Rational Inattention Foundation} The parameter $\varepsilon$ can be motivated using the rational inattention framework of \citet{sims2003implications}. The key insight is that the two forecasting rules differ not only in their accuracy but also in their cognitive demands. The CB-anchored rule requires recalling a single fixed number: the announced target. The backward-looking rule requires tracking realized inflation, a time-varying process that demands continuous updating. When forecast accuracy is similar, a cognitively constrained agent will prefer the simpler rule. In the rational inattention framework, agents face a shadow cost $\lambda$ per unit of information processed, measured in the same units as the loss function. An agent will prefer the CB-anchored rule whenever:
  
\begin{equation}
L^{CB}_t - L^{BL}_t < \lambda \cdot \Delta H
\end{equation}

where $\Delta H$ is the additional information (in bits) required to track inflation rather than recall a constant. This yields our tie-breaking parameter: $\varepsilon = \lambda \cdot \Delta H$.

\paragraph{Calibration} We now show that $\varepsilon = 10^{-4}$ is consistent with information costs estimated in the rational inattention literature. Throughout, inflation is measured in decimal form (so $\pi = 0.02$ represents 2\% inflation), and losses are in squared decimal units.

First, note that $\varepsilon = 10^{-4} = (0.01)^2$. Since inflation is measured in decimals, this corresponds to the squared forecast error from a one-percentage-point miss --- the CB-anchored rule wins ties as long as it is not more than one percentage point less accurate than the backward-looking alternative. Concretely: suppose the target is 2\% and last period's inflation was 3\%. Without the simplicity premium, agents would be indifferent between rules when inflation comes in at 2.5\% (the midpoint). With $\varepsilon = 10^{-4}$, the threshold shifts to 3\%: agents stick with the CB-anchored rule unless realized inflation matches or exceeds last period's level, validating the backward-looking forecast.

Second, this magnitude aligns with estimates from the rational inattention literature. \citet{sims2003implications} calibrates models where agents with moderate information-processing capacity (0.3--3 bits per period) optimally tolerate 10--50\% of variance unexplained. With U.S.\ inflation variance approximately $\text{Var}(\pi) \approx (0.02)^2 = 4 \times 10^{-4}$, tolerating 30\% unexplained implies:

\begin{equation}
\text{Information cost} \approx 0.3 \times 4 \times 10^{-4} = 1.2 \times 10^{-4}
\end{equation}

This is the same order of magnitude as our $\varepsilon = 10^{-4}$.

Third, we can express this in bits. The entropy rate of tracking a Gaussian process with innovation standard deviation $\sigma_\varepsilon \approx 0.01$ (one percentage point) is approximately 2--3 bits per period. If agents tolerate information costs equivalent to $(0.01)^2$ per 2--3 bits, the implied shadow cost is:

\begin{equation}
\lambda \approx \frac{10^{-4}}{2.5 \text{ bits}} \approx 4 \times 10^{-5} \text{ per bit}
\end{equation}

In sum, $\varepsilon = 10^{-4}$ represents an information cost consistent with agents who avoid cognitive effort when the stakes are small, but switch to more demanding strategies when simpler rules fail persistently.

%%%%%%%%%%%%%%%%%%%%%%%%%%%%%%%%%%%%%%%%%%%%%%%%%%%%%%%%%%%%%%%%%%%%%%%%%%%%%%%%%%%%%%%%%%%%%%%%%%%%%%%%%%%%%%%%%%%%%%%%%%%%%%%%%%%%%%%%%%%%%%%%%%%%%%%%%%%%%%%%%%%%%%%%%%%%%%%%%%%%%%%%%%%%%%%%%%%%%%%%%%%%%%%%%%%%%%%%%%%%%%%%%%%%%%%%%%%%%%%%%%%%%%%%%%%%%%%%%%%%%%%%%%%%%%%%%%%%%%%%%%%%%%%%%%%%%%%%%%%%%%%%%%%%%%%%%%%%%%%%%%%%%%%%%%%%%%%%%%%%%%%%%%%%%%%%%%%%%%%%%%%%%%%%%%%%%%%%%%%%%%%%%%%%%%%%%%%%%%%%%%%%%%%%%%%%%%%%%%%%%%%%%%%%%%%%%%%%%%%%%%%%%%%%%%%%%%%%%%%%%%%%%%%%%%%%%%%%%%%%%%%%%%%%%%%%%%%%%%%%%%%%%%%%%%%%%%%%%%%%%%%%%%%%%%%%%%%%%%%%%%%%%%%%%%%%%%%%%%%%%%%%%%%%%%%%%%%%%%%%%%%%%%%%%%%%%%%%%%%%%%%%%%%%%%%%%%%%%%%%%%%%%%%%%%%%%%%%%%%%%%%%%%%%%%%%%%%%%%%%%%%%%%%%%%%%%%%%%%%%%%%%%%%%%%%%%%%%%%%%%%%%%%%%%%%%%%%%%%%%%%%%%%%%%%%%%%%%%%%%%%%%%%%%%%%%%%%

\section{Numerical Solution Methods}
\label{app:computational}

This appendix documents the computational algorithm used to solve the model when FIRE agents are ``sophisticated'': that is, when they observe the credibility state $\theta_t$ and form expectations taking its evolution into account. We first explain why this creates a fixed-point problem, then describe our solution algorithm.

\paragraph{The Fixed-Point Problem} In the baseline model, FIRE agents are ``naive'': they solve the standard New Keynesian model assuming $\lambda = 1$ (all agents rational), ignoring the bounded rationality population and the dynamics of $\theta_t$. This simplification is common in the heterogeneous expectations literature \citep[e.g.,][]{degrauwe2011}  and yields a tractable recursive structure. With sophisticated FIRE, agents recognize that aggregate expectations depend on $\theta_t$:
  
\begin{equation}
\mathbb{E}_t[\pi_{t+1}] = \lambda \cdot \mathbb{E}^{FIRE}_t[\pi_{t+1}] + (1-\lambda) \cdot \left[\theta_t \cdot \pi^* + (1-\theta_t) \cdot \pi_{t-1}\right]
\end{equation}

Since $\theta_t$ evolves based on realized inflation, and inflation depends on expectations, this creates a fixed-point problem: FIRE expectations depend on the path of $\theta$, which depends on inflation outcomes via MAB learning, which in turn depends on aggregate expectations including FIRE. We solve this via iteration.

\paragraph{Solution Algorithm} The algorithm uses the naive model solution as a warm start, then iterates until convergence. This warm start is effective because the naive solution is already close to the sophisticated solution in most states; the key difference arises during de-anchoring episodes, where sophisticated FIRE agents anticipate credibility erosion and adjust expectations preemptively.

\textit{Step 0: Warm Start.} Solve the naive FIRE model to obtain initial guesses $\{\mathbb{E}^{FIRE,0}_t\}_{t=1}^T$.

\textit{Step 1: Iterate.} For iteration $n = 1, 2, \ldots$:

\begin{enumerate}
\item[(a)] \textbf{Forward simulation.} Given current FIRE expectations $\{\mathbb{E}^{FIRE,n-1}_t\}$, simulate forward: for each $t = 1, \ldots, T$, compute aggregate expectations, solve the NK system for $(\pi_t, y_t)$, and update $\theta_t$ via MAB learning. Store the resulting paths $\{\pi^n_t, \theta^n_t\}$.
\item[(b)] \textbf{Backward solution.} Given $\{\pi^n_t, \theta^n_t\}$ from forward simulation, solve backward: set terminal condition $\mathbb{E}^{FIRE}_{T-1}[\pi_T] = \pi^*$, then for $t = T-2, \ldots, 1$, compute $\pi_{t+1}$ from the NK system given $\theta^n_{t+1}$ and $\mathbb{E}^{FIRE,n}_{t+1}$, and set $\mathbb{E}^{FIRE,n}_t[\pi_{t+1}] = \pi_{t+1}$.
\item[(c)] \textbf{Convergence and dampening.} Compute $\delta^n = \max_t |\mathbb{E}^{FIRE,n}_t - \mathbb{E}^{FIRE,n-1}_t|$. If $\delta^n < \text{tol}$, stop. Otherwise, apply dampening: $\mathbb{E}^{FIRE,n}_t \leftarrow \alpha \cdot \mathbb{E}^{FIRE,n}_t + (1-\alpha) \cdot \mathbb{E}^{FIRE,n-1}_t$.
\end{enumerate}

\paragraph{Implementation Details} We use dampening parameter $\alpha = 0.3$ for numerical stability, convergence tolerance $\text{tol} = 10^{-8}$, and maximum iterations $N_{\max} = 200$. In practice, convergence typically occurs within 40--100 iterations. Each iteration requires two passes through the simulation horizon (one forward, one backward), each involving a $2 \times 2$ linear system at each period. For $T = 120$ quarters (30 years), total runtime is under one second on a standard laptop, allowing extensive Monte Carlo analysis.

Sophisticated FIRE amplifies all de-anchoring dynamics documented in the main text: credibility troughs are deeper, recovery takes longer, and inflation peaks are higher. All qualitative results are preserved: the mechanism operates identically, but magnitudes are larger.

\paragraph{Softmax Smoothing for Policy Optimization} The optimal policy analysis in Section~\ref{sec:applications} searches over the policy coefficient $\phi_\pi$ to minimize welfare loss. This optimization is complicated by the discrete MAB switching rule: the hard indicator $\mathds{1}[L^{CB}_t \leq L^{BL}_t + \varepsilon]$ creates discontinuities in the mapping from policy parameters to welfare outcomes. A small change in $\phi_\pi$ can flip the indicator in some period, causing a discrete jump in the welfare function. The resulting loss surface is non-smooth, producing spurious non-monotonicity in the optimal policy curve.

We address this by replacing the hard indicator with a softmax switching probability:

\begin{equation}
P(CB) = \frac{1}{1 + \exp(\gamma (L^{CB}_t - L^{BL}_t - \varepsilon))}
\end{equation}

where $\gamma > 0$ is an inverse temperature parameter controlling the sharpness of the switch, and updating credibility as $\theta_{t+1} = (1-\eta)\theta_t + \eta \cdot P(CB)$. This renders the welfare function continuously differentiable in $\phi_\pi$.

We set $\gamma = 25{,}000$ to balance two objectives. First, it preserves model dynamics: the transition from $P(CB) \approx 0.01$ to $P(CB) \approx 0.99$ occurs over a loss difference of approximately $\pm 0.0002$, effectively a sharp switch for typical loss magnitudes. Second, it ensures smooth optimization: the sigmoid is continuously differentiable, so grid search identifies a unique minimum at each credibility level. 

%%%%%%%%%%%%%%%%%%%%%%%%%%%%%%%%%%%%%%%%%%%%%%%%%%%%%%%%%%%%%%%%%%%%%%%%%%%%%%%%%%%%%%%%%%%%%%%%%%%%%%%%%%%%%%%%%%%%%%%%%%%%%%%%%%%%%%%%%%%%%%%%%%%%%%%%%%%%%%%%%%%%%%%%%%%%%%%%%%%%%%%%%%%%%%%%%%%%%%%%%%%%%%%%%%%%%%%%%%%%%%%%%%%%%%%%%%%%%%%%%%%%%%%%%%%%%%%%%%%%%%%%%%%%%%%%%%%%%%%%%%%%%%%%%%%%%%%%%%%%%%%%%%%%%%%%%%%%%%%%%%%%%%%%%%%%%%%%%%%%%%%%%%%%%%%%%%%%%%%%%%%%%%%%%%%%%%%%%%%%%%%%%%%%%%%%%%%%%%%%%%%%%%%%%%%%%%%%%%%%%%%%%%%%%%%%%%%%%%%%%%%%%%%%%%%%%%%%%%%%%%%%%%%%%%%%%%%%%%%%%%%%%%%%%%%%%%%%%%%%%%%%%%%%%%%%%%%%%%%%%%%%%%%%%%%%%%%%%%%%%%%%%%%%%%%%%%%%%%%%%%%%%%%%%%%%%%%%%%%%%%%%%%%%%%%%%%%%%%%%%%%%%%%%%%%%%%%%%%%%%%%%%%%%%%%%%%%%%%%%%%%%%%%%%%%%%%%%%%%%%%%%%%%%%%%%%%%%%%%%%%%%%%%%%%%%%%%%%%%%%%%%%%%%%%%%%%%%%%%%%%%%%%%%%%%%%%%%%%%%%%%%%%%%%%%%%%%

\section{Robustness Analysis}
\label{app:robustness}

This appendix verifies that the time-varying persistence estimates in the main text are robust to alternative estimation methods and window size choices.

\paragraph{State-Space Approach} The main text estimates time-varying inflation persistence using rolling 40-quarter AR(1) regressions. While intuitive, this approach requires an arbitrary window choice and provides no formal uncertainty quantification. As a robustness check, we re-estimate persistence using a state-space model with Kalman filtering, which lets the data determine how much persistence varies over time.

We specify a time-varying parameter AR(1) model:

\begin{align}
\pi_t &= c + \rho_t \pi_{t-1} + \varepsilon_t, \quad \varepsilon_t \sim N(0, \sigma^2_{\varepsilon}) \\
\rho_t &= \rho_{t-1} + \omega_t, \quad \omega_t \sim N(0, \sigma^2_{\omega})
\end{align}

The observation equation is a standard AR(1) with time-varying coefficient $\rho_t$; the state equation specifies that persistence follows a random walk. We estimate the variance parameters $(\sigma^2_{\varepsilon}, \sigma^2_{\omega})$ via maximum likelihood, then apply the Kalman filter (forward pass) and Rauch-Tung-Striebel smoother (backward pass) to obtain smoothed estimates $\hat{\rho}_{t|T}$ that use all available information. The ratio $\sigma^2_{\omega}/\sigma^2_{\varepsilon}$ governs the smoothness of the estimated path; we estimate this from the data rather than imposing it.

\paragraph{Results} Maximum likelihood yields $\hat{\sigma}^2_{\varepsilon} = 5.12$ and $\hat{\sigma}^2_{\omega} = 0.0036$, implying a low signal-to-noise ratio (0.0007) consistent with gradual credibility transitions. Table~\ref{tab:kalman_comparison} compares the Kalman estimates with the rolling window estimates from the main text.

\begin{table}[htbp]
\centering
\caption{Comparison: Kalman TVP vs.\ Rolling Window Persistence}
\label{tab:kalman_comparison}
\begin{tabular}{lccc}
\toprule
Period & Kalman $\bar{\rho}$ & Rolling $\bar{\rho}$ & Difference \\
\midrule
Great Inflation (1968--1979) & 0.703 & 0.729 & $-$0.026 \\
Volcker Era (1980--1984) & 0.490 & 0.671 & $-$0.181 \\
Great Moderation (1985--2007) & 0.145 & 0.255 & $-$0.109 \\
Post-Crisis (2008--2019) & 0.009 & $-$0.045 & $+$0.054 \\
COVID Era (2020--2024) & 0.352 & 0.392 & $-$0.040 \\
\midrule
Correlation & \multicolumn{3}{c}{0.86} \\
\bottomrule
\end{tabular}
\end{table}

The correlation between methods is 0.86, indicating strong agreement on the overall pattern. Both approaches identify high persistence during the Great Inflation, declining persistence through the 1980s and 1990s, a trough during the anchored-expectations era, and rising persistence during the COVID episode. The Kalman estimates are systematically smoother, particularly during the Volcker transition, reflecting optimal weighting of signal versus noise.

\paragraph{Window Size Sensitivity} As a complementary check, we verify that the rolling window estimates are stable across different window sizes. Table~\ref{tab:window_sensitivity} reports results using 32, 40, and 48 quarter windows.

\begin{table}[htbp]
\centering
\caption{Rolling Persistence Estimates by Window Size}
\label{tab:window_sensitivity}
\begin{tabular}{lcccc}
\toprule
Window (quarters) & Mean $\hat{\rho}$ & Std.\ Dev. & Min & Max \\
\midrule
32 & 0.268 & 0.331 & $-$0.413 & 0.867 \\
40 & 0.290 & 0.344 & $-$0.354 & 0.859 \\
48 & 0.296 & 0.349 & $-$0.358 & 0.867 \\
\bottomrule
\end{tabular}
\end{table}

Mean persistence estimates are stable across window sizes, ranging from 0.268 to 0.296. Combined with the Kalman filter validation, these results provide evidence that the main text conclusions about time-varying persistence are not artifacts of methodological choices.

%%%%%%%%%%%%%%%%%%%%%%%%%%%%%%%%%%%%%%%%%%%%%%%%%%%%%%%%%%%%%%%%%%%%%%%%%%%%%%%%%%%%%%%%%%%%%%%%%%%%%%%%%%%%%%%%%%%%%%%%%%%%%%%%%%%%%%%%%%%%%%%%%%%%%%%%%%%%%%%%%%%%%%%%%%%%%%%%%%%%%%%%%%%%%%%%%%%%%%%%%%%%%%%%%%%%%%%%%%%%%%%%%%%%%%%%%%%%%%%%%%%%%%%%%%%%%%%%%%%%%%%%%%%%%%%%%%%%%%%%%%%%%%%%%%%%%%%%%%%%%%%%%%%%%%%%%%%%%%%%%%%%%%%%%%%%%%%%%%%%%%%%%%%%%%%%%%%%%%%%%%%%%%%%%%%%%%%%%%%%%%%%%%%%%%%%%%%%%%%%%%%%%%%%%%%%%%%%%%%%%%%%%%%%%%%%%%%%%%%%%%%%%%%%%%%%%%%%%%%%%%%%%%%%%%%%%%%%%%%%%%%%%%%%%%%%%%%%%%%%%%%%%%%%%%%%%%%%%%%%%%%%%%%%%%%%%%%%%%%%%%%%%%%%%%%%%%%%%%%%%%%%%%%%%%%%%%%%%%%%%%%%%%%%%%%%%%%%%%%%%%%%%%%%%%%%%%%%%%%%%%%%%%%%%%%%%%%%%%%%%%%%%%%%%%%%%%%%%%%%%%%%%%%%%%%%%%%%%%%%%%%%%%%%%%%%%%%%%%%%%%%%%%%%%%%%%%%%%%%%%%%%%%%%%%%%%%%%%%%%%%%%%%%%%%%%%%%


\section{Bayesian Learning: An Alternative Microfoundation}
\label{app:bayesian}

This appendix presents an alternative expectation formation mechanism based on Bayesian learning. While the main text uses a multi-armed bandit (MAB) framework where agents discretely choose between forecasting rules, here agents form weighted combinations of information sources and update weights via Bayes' rule. The connection between MAB and Bayesian learning is deep: Thompson sampling, one of the most successful bandit algorithms, is itself a Bayesian method \citep{thompson1933, agrawal2012}. By showing that our main results survive under this more conventional microfoundation, we find that the key findings depend on the economic logic of performance-based rule evaluation, not on the specific assumptions of the MAB framework.

\paragraph{Setup} bounded rationality agents face a signal extraction problem. They observe two information sources about future inflation: the central bank's announced target $\pi^*$, and the most recently realized inflation $\pi_{t-1}$. Unlike the MAB framework where agents discretely choose one rule, here agents form expectations as a weighted average:

\begin{equation}
\mathbb{E}^{BR}_t[\pi_{t+1}] = w \cdot \pi^* + (1-w) \cdot \pi_{t-1}
\end{equation}

where $w \in [0,1]$ is the weight on the central bank target. When $w = 1$, the agent is fully anchored; when $w = 0$, purely backward-looking.

We model beliefs about the appropriate weight using a Beta distribution: $w \sim \text{Beta}(\alpha_t, \beta_t)$. The Beta distribution is natural here: it is defined on $[0,1]$, matching the domain of $w$, and is conjugate to the Bernoulli likelihood, yielding closed-form updates. The parameters have intuitive interpretations: $\alpha_t$ represents accumulated evidence favoring the central bank target, while $\beta_t$ represents evidence favoring backward-looking forecasts. The posterior mean,
  
\begin{equation}
\bar{w}_t = \frac{\alpha_t}{\alpha_t + \beta_t}
\end{equation}

is analogous to $\theta_t$ in the MAB model.

\paragraph{Updating Rule} Each period, agents observe which forecasting rule was more accurate. Define forecast errors $L^{CB}_t = (\pi_t - \pi^*)^2$ and $L^{BL}_t = (\pi_t - \pi_{t-1})^2$, matching the MAB specification. Let $s_t = \mathbf{1}\{L^{CB}_t \leq L^{BL}_t + \varepsilon\}$ indicate whether the central bank's forecast was at least as accurate (with the same simplicity bias $\varepsilon$ as the MAB specification). The posterior parameters evolve as:

\begin{align}
\alpha_{t+1} &= \alpha_t + s_t \\
\beta_{t+1} &= \beta_t + (1 - s_t)
\end{align}

When the CB-anchored rule is more accurate, $\alpha$ increases; when backward-looking extrapolation performs better, $\beta$ increases. Learning speed emerges endogenously from the precision of prior beliefs: diffuse priors ($\alpha_0 + \beta_0$ small) produce fast learning; precise priors produce gradual updating.

\paragraph{Comparison with MAB} To compare the frameworks, we simulate both models using identical shock sequences and calibration, with initial conditions chosen so that $\bar{w}_0 = \theta_0$. Figure~\ref{fig:mab_vs_bayesian} presents the results across three initial credibility levels.

\begin{figure}[htbp]
\centering
\includegraphics[width=\textwidth]{../output/figures/appendix_E/figure_E1_mab_vs_bayesian.png}
\caption{MAB vs Bayesian Learning: Comparison of Aggregate Dynamics. Each row corresponds to a different initial credibility level. Column (a): credibility state ($\theta_t$ for MAB, $\bar{w}_t$ for Bayesian). Column (b): inflation. Column (c): output gap. Both models face a 2\% cost-push shock at $t=5$.}
 \label{fig:mab_vs_bayesian}
\end{figure}

The two frameworks generate remarkably similar aggregate dynamics. Across all initial conditions, the correlation between $\theta_t$ and $\bar{w}_t$ exceeds 0.87, and the root mean squared error between inflation paths is below 0.25 percentage points. Both models exhibit the same qualitative patterns: credibility drops sharply after the shock as backward-looking forecasts temporarily dominate, then gradually recovers as inflation returns to target. The key asymmetry (rapid de-anchoring followed by slow re-anchoring) emerges naturally in both frameworks. This suggests that the credibility dynamics documented in the main text are robust to the choice of learning mechanism: whether agents use simple performance tracking or Bayesian updating, the same qualitative patterns emerge.

\paragraph{Discussion} Given the similarity in aggregate dynamics, why favor MAB in the main text? Both approaches have merits. The Bayesian framework is more familiar to economists and offers a natural measure of uncertainty through the posterior variance. The MAB framework, however, has advantages for our purposes. First, it is simpler: $\theta_t$ has a direct interpretation as the fraction of agents using each rule. Second, discrete switching between rules (rather than continuous weight adjustment) better matches experimental evidence that subjects adopt heuristics wholesale rather than blend them continuously \citep{hommes2021}. Third, the sharp transitions generated by MAB are more consistent with rapid de-anchoring observed in historical episodes. The near-equivalence of aggregate outcomes under both frameworks suggests our main results are robust to this modeling choice.

\end{document}